\RequirePackage[thinlines]{easybmat} %-- muss aufgrund eines Bugs vor etex und tikz geladen werden
\documentclass[a4paper, twoside, headsepline, index=totoc,toc=listof, fontsize=10pt, cleardoublepage=empty, headinclude, DIV=13, BCOR=13mm, titlepage,]{scrartcl}
\usepackage{scrtime} % KOMA, Uhrzeit ermoeglicht
\usepackage{scrpage2} % wie fancyhdr, nur optimiert auf KOMA-Skript, leicht andere Syntax
\usepackage{etoolbox}
\usepackage{letltxmacro}

% Compiler abhaengige IFs
\usepackage{ifxetex,ifluatex}
\newif\ifxetexorluatex
\ifxetex
  \xetexorluatextrue
\else
  \ifluatex
    \xetexorluatextrue
  \else
    \xetexorluatexfalse
  \fi
\fi

%--Farbdefinitionen
%-- muss vor tikz geladen werden
\usepackage[usenames, table, x11names]{xcolor}
\definecolor{dark_gray}{gray}{0.45}
\definecolor{light_gray}{gray}{0.6}

%--Zum Zeichnen
%-- muss vor polyglossia bzw. babel geladen werden
\usepackage{tikz}
\usepackage{tikz-cd}
\usetikzlibrary{external}
\tikzset{>=latex}
\usetikzlibrary{shapes,arrows,intersections}
\usetikzlibrary{calc,3d}
\usetikzlibrary{decorations.pathreplacing,decorations.markings}
\usetikzlibrary{angles}

%-- Konfiguration von tikzexternalize
\tikzexternalize[prefix=tikz/, up to date check=diff]

\ifxetexorluatex
\tikzset{external/system call={xelatex \tikzexternalcheckshellescape %-- verwende LuaLaTeX, wegen dynamischer Speicherallokation
    -halt-on-error -interaction=batchmode -jobname "\image" "\texsource"}}
\else
\tikzset{external/system call={pdflatex \tikzexternalcheckshellescape 
    -halt-on-error -interaction=batchmode -jobname "\image" "\texsource"}}
\fi

%-- tikzexternalize fuer tikzcd deaktivieren, da inkompatibel
\AtBeginEnvironment{tikzcd}{\tikzexternaldisable}
\AtEndEnvironment{tikzcd}{\tikzexternalenable}

%-- um Inkompatibilitaeten von quotes und polyglossia bzw. babel zu vermeiden
\tikzset{
  every picture/.append style={
    execute at begin picture={\shorthandoff{"}},
    execute at end picture={\shorthandon{"}}
  }
}
\usetikzlibrary{quotes}
\usepackage{pgfplots}
\usepgfplotslibrary{colormaps}
\pgfplotsset{compat=1.10}



%-- Mathe
\usepackage{mathtools} % beinhaltet amsmath
\mathtoolsset{showonlyrefs, centercolon}
\newtagform{brackets}[\textbf]{[}{]}
\usetagform{brackets}
\usepackage{wasysym}
\usepackage{amssymb} %zusätzliche Symbole
\usepackage{latexsym} %zusätzliche Symbole
\usepackage{stmaryrd} %für Blitz
\usepackage{nicefrac} %schräge Brüche
\usepackage{cancel} %Befehle zum Durchstreichen
\usepackage{mathdots}
% \usepackage{bm}
\DeclareSymbolFont{bbold}{U}{bbold}{m}{n}
\DeclareSymbolFontAlphabet{\mathbbold}{bbold}
\newcommand{\mathds}[1]{\mathbb{#1}} %Um Kompatibilitaet mit frueheren Benutzung von dsfont herzustellen
\newcommand{\ind}{\mathbbold{1}} % charakteristische-Funktion-Eins
\def\mathul#1#2{\color{#1}\underline{{\color{black}#2}}\color{black}} %farbiges Untersteichen im Mathe-Modus

%-- Alles was mit Schrift und XeteX zu tun hat
\ifxetexorluatex
    % XeLaTeX or LuaTeX
	\ifxetex
		\usepackage{mathspec} %beinhaltet fontspec 
	\else
		\usepackage[no-math]{fontspec}
	\fi
    \usepackage{polyglossia} %babel-ersatz
    \setmainlanguage[spelling=new,babelshorthands=false]{german}
	\newcommand\glqq{"}
	\newcommand\grqq{"}
	\defaultfontfeatures{Mapping=tex-text, WordSpace={1.4}} %
    \setmainfont[Ligatures=Common, BoldFont={* Bold}, ItalicFont={* Light Italic},SmallCapsFont={Linux Libertine O}, SmallCapsFeatures={Letters=SmallCaps}]{Source Sans Pro}
	\setsansfont[Scale=MatchLowercase,Ligatures=Common, BoldFont={* Medium}]{Ubuntu}
	\ifxetex
		\setallmonofonts[Scale=MatchLowercase, ItalicFont={*}]{Consolas} 
	\else
		\setmonofont[Scale=MatchLowercase, ItalicFont={*}]{Consolas}
	\fi
	\usepackage{xltxtra}
	\usepackage{fontawesome}
	\usepackage{microtype}
\else
    % default: pdfLaTeX
    \usepackage[ngerman]{babel}
    \usepackage[T1]{fontenc}
    \usepackage[utf8]{inputenc}
	\usepackage{textcomp} %verhindert ein paar Fehler bei den Fonts
    \usepackage[babel=true, tracking=true]{microtype}
	\usepackage{lmodern}
	\renewcommand{\familydefault}{\sfdefault} 
\fi


%--Mixed
\usepackage[neverdecrease]{paralist}
\usepackage[german=quotes]{csquotes}
\usepackage{makeidx}
\usepackage{booktabs}
\usepackage{wrapfig}
\usepackage{float}
\usepackage[margin=10pt, font=small, labelfont={sf, bf}, format=plain, indention=1em]{caption}
\captionsetup[wrapfigure]{name=Abb. }
\usepackage{stackrel}
\usepackage{ifthen}
\usepackage{multicol}
\usepackage{qrcode}
\flushbottom

%--Hyphenation
\hyphenation{di-men-si-o-na-le}


%--Unterstreichung
\usepackage[normalem]{ulem}
\setlength{\ULdepth}{1.8pt}

%--Indexverarbeitung
\newcommand{\bet}[1]{\textbf{#1}}
\newcommand{\Index}[1]{\textbf{#1}\index{#1}}
\makeindex
\setindexpreamble{{\noindent \itshape Die \emph{Seitenzahlen} sind mit Hyperlinks zu den entsprechenden Seiten versehen, also anklickbar} \ifxetexorluatex \faHandUp\fi \par \bigskip}
\renewcommand{\indexpagestyle}{scrheadings}


%--Marginnote & todonotes
\usepackage{marginnote}
\renewcommand*{\marginfont}{\color{dark_gray} \itshape \footnotesize }
\usepackage[textsize=small]{todonotes}

%--Todonotes mit tikz/externalize kompatibel machen
\LetLtxMacro{\oldtodo}{\todo}
\renewcommand{\todo}[2][]{\tikzexternaldisable\oldtodo[#1]{#2}\tikzexternalenable}
\LetLtxMacro{\oldmissingfigure}{\missingfigure}
\renewcommand{\missingfigure}[2][]{\tikzexternaldisable\oldmissingfigure[{#1}]{#2}\tikzexternalenable}

%--Konfiguration von Hyperref pdfstartview=FitH, 
\usepackage[hidelinks, pdfpagelabels,  bookmarksopen=true, bookmarksnumbered=true, linkcolor=black, urlcolor=SkyBlue2, plainpages=false,pagebackref, citecolor=black, hypertexnames=true, pdfauthor={Jannes Bantje}, pdfborderstyle={/S/U}, linkbordercolor=SkyBlue2, colorlinks=false]{hyperref}

% \let\orighref\href
\ifxetexorluatex
\newcommand{\hrefsym}[2]{\href{#1}{\texttt{#2\,\faExternalLink}}}
\newcommand{\hrefsymmail}[2]{\href{#1}{\texttt{\faEnvelope\,#2}}}
\renewcommand{\url}[1]{\hrefsym{#1}{\nolinkurl{#1}}}
\fi
\providecommand{\hrefsym}{\href}
% \usepackage{cleveref}
% \crefformat{footnote}{#2\footnotemark[#1]#3}

%--Römische Zahlen
\newcommand{\RM}[1]{\MakeUppercase{\romannumeral #1{}}}

%--Punkte (nach hyperref)
\usepackage{ellipsis}

%%--Abkürzungen etc.
\newcommand{\light}{\text{\Large $\lightning$}}

%-- Definitionen von weiteren Mathe-Befehlen
\DeclareMathOperator{\re}{Re} %Realteil
\DeclareMathOperator{\im}{Im} %Imaginaerteil
\DeclareMathOperator{\id}{id} %identische Abbildung
\DeclareMathOperator{\Sp}{Sp} %Spur
\DeclareMathOperator{\sgn}{sgn} %Signum
\DeclareMathOperator{\alt}{Alt} %Alternierende n-linearFormen
\DeclareMathOperator{\End}{End} %Endomorphismen
\DeclareMathOperator{\Vol}{Vol} %Volumen
\DeclareMathOperator{\dom}{dom} %Domain
\DeclareMathOperator{\Hom}{Hom} %Homomorphismen
\DeclareMathOperator{\bild}{Bild} %Bild
\DeclareMathOperator{\Kern}{Kern }
\DeclareMathOperator{\rg}{rg} %Rang
\DeclareMathOperator{\Rg}{Rg} %Rang
\DeclareMathOperator{\diam}{diam} %Durchmesser
\DeclareMathOperator{\dist}{dist} %Distanz
\DeclareMathOperator{\grad}{grad} %Gradient
\DeclareMathOperator{\dive}{div} %Gradient
\DeclareMathOperator{\rot}{rot} %Rotation
\DeclareMathOperator{\hess}{Hess} %Hesse-Matrix
\DeclareMathOperator{\koker}{Koker} %Kokern
\DeclareMathOperator{\aut}{Aut} %Automorphismen
\DeclareMathOperator{\ord}{ord} %Ordnung
\DeclareMathOperator{\ggT}{ggT} %ggT
\DeclareMathOperator{\kgV}{kgV} %kgV
\DeclareMathOperator{\Gr}{Gr} %Gerade
\DeclareMathOperator{\Kr}{Kr} %Kreis
\DeclareMathOperator{\Char}{char} %Charakteristik
\DeclareMathOperator{\Aut}{Aut} %Automorphismen
% \DeclareMathOperator{\D}{D} %formale Ableitung
\newcommand{\D}{\ensuremath{\mathrm{D}\mkern-1.5mu}}
\newcommand{\mathd}{\ensuremath{\mathrm{d}\mkern-1.0mu}}
\newcommand{\Tmap}{\ensuremath{\mathrm{T}\mkern-1.0mu}}
\newcommand{\opL}{\ensuremath{\mathrm{L}\mkern-0.6mu}}
\DeclareMathOperator{\cov}{cov} %Kovarianz
\DeclareMathOperator{\Gal}{Gal} %Galois
\DeclareMathOperator{\supp}{supp} %Träger
\DeclareMathOperator{\Sym}{Sym} %Symmetrische Gruppe
\DeclareMathOperator{\Zyl}{Zyl} %Zylinder
\DeclareMathOperator{\Mod}{Mod} %Moduln
\DeclareMathOperator{\EPK}{EPK} %Einpunktkompaktifizierung
\DeclareMathOperator{\conj}{conj} %Konjugation
\DeclareMathOperator{\ann}{ann} %Annulator
\DeclareMathOperator{\tor}{tor} %Torsionsmodul
\DeclareMathOperator{\Int}{Int} %Interior/Inneres
\DeclareMathOperator{\Br}{Br} %Brauerergruppe
\DeclareMathOperator{\GL}{GL} %allgemeine lineare Gruppe
\DeclareMathOperator{\SL}{SL} %spezielle lineare Gruppe
\DeclareMathOperator{\SU}{SU} %spezielle unitäre Gruppe
\DeclareMathOperator{\Tr}{Tr} %Spur/Trace
\DeclareMathOperator{\conv}{conv} %Konvexe Teilmenge
% \DeclareMathOperator{\tr}{tr}
\newcommand{\tr}{\mathrm{tr}}
\newcommand{\w}{\mathrm{w}}
\newcommand{\weakT}[1]{\ensuremath{\mathcal{T}_{#1}^{\mkern+1.0mu\text{\raisebox{0.4ex}{$\mathrm{w}$}}}}}
\newcommand{\weakTstar}[1]{\ensuremath{\mathcal{T}_{#1}^{\mkern+1.0mu\text{\raisebox{0.4ex}{$\mathrm{w}$}}^*}}}
\newcommand{\rand}[1]{\ensuremath{\partial^{\scriptscriptstyle #1}}}


% -- Kategorien
\DeclareMathOperator{\Ob}{Ob} %Objekte
\DeclareMathOperator{\Mor}{Mor} %Morphismen

% \DeclareMathOperator{\TOP}{TOP} %Topologische Räume
\newcommand{\TOP}{\textsc{Top}}
% \DeclareMathOperator{\HTOP}{HTOP} %Topologische Räume + Homotopieäquivalenz
\newcommand{\HTOP}{\textsc{HTop}}
% \DeclareMathOperator{\VR}{VR} %Vektorräume
\newcommand{\VR}{\textsc{VR}}
% \DeclareMathOperator{\MOD}{MOD} %Moduln
\newcommand{\MOD}{\textsc{Mod}}
% \DeclareMathOperator{\MENGEN}{MENGEN} %Mengen
\newcommand{\MENGEN}{\textsc{Mengen}}
% \DeclareMathOperator{\MAN}{MAN} %Mannigfaltigkeiten
\newcommand{\MAN}{\textsc{Man}}
% \DeclareMathOperator{\GRUPPEN}{GRUPPEN} %Gruppen
\newcommand{\GRUPPEN}{\textsc{Gruppen}}
% \DeclareMathOperator{\ABELGRUPPEN}{ABEL.GRUPPEN} %abelsche Gruppen
\newcommand{\ABELGRUPPEN}{\textsc{Abel.Gruppen}}
% \DeclareMathOperator{\KAT}{KAT} %Kategorien
\newcommand{\KAT}{\textsc{Kat}}
% \DeclareMathOperator{\FUN}{FUN} %Funktoren
\newcommand{\FUN}{\textsc{Fun}}


\newcommand{\op}{\mathrm{op}}

\newcommand{\Ker}{\ensuremath{\Kern \,}}
\newcommand{\zirkel}{\rotatebox[origin = c]{-90}{$\varangle$}}
\newcommand{\wirkung}{\!\curvearrowright\!}

\newcommand{\oE}{\OE}

%--Beweisende
\newcommand{\bewende}{\ifmmode \tag*{$\square$} \else \hfill $\square$ \fi}

%--Volumen
\newcommand{\vol}[1]{\ensuremath{\Vol_{#1}}}

%--Skalarprodukt
\DeclarePairedDelimiterX\sprod[2]{\langle}{\rangle}{#1\,\delimsize\vert\,#2}

%--Betrag, Gaußklammer
\DeclarePairedDelimiter{\abs}{\lvert}{\rvert}
\DeclarePairedDelimiter{\ceil}{\lceil}{\rceil}

%--Norm
\DeclarePairedDelimiter\doppelstrich{\Vert}{\Vert}
\newcommand{\norm}[2][\relax]{
\ifx#1\relax \ensuremath{\doppelstrich*{#2}}
\else \ensuremath{\doppelstrich*{#2}_{#1}}
\fi}


%--Umklammern
\DeclarePairedDelimiter\enbrace{(}{)}
\DeclarePairedDelimiter\benbrace{[}{]}
\DeclarePairedDelimiter\lenbrace{<}{>}

%--Mengen
\DeclarePairedDelimiterX\mengenA[1]{\lbrace}{\rbrace}{#1}
\DeclarePairedDelimiterX\mengenB[2]{\lbrace}{\rbrace}{#1\, \delimsize\vert \, #2}

\makeatletter
\newcommand{\set}[2][\relax]{
\ifx#1\relax \ensuremath{
\mengenA*{#2}}
\else \ensuremath{%
  \mengenB*{#1}{#2}}
\fi}
\makeatother


%--offen und abgeschlossen
\newcommand{\off}{\! \stackrel[\text{offen}]{}{\subset} \!}
\newcommand{\abg}{\! \stackrel[\text{abg.}]{}{\subset} \!}

%--Differential
\newcommand{\diff}[2]{\ensuremath{\frac{{\partial #1}}{{\partial #2}} }}
\newcommand{\diffd}[2]{\ensuremath{\frac{\mathd #1}{\mathd #2} }}

%--Diagonale Linien für Matrizen
\newcommand\tikzmark[1]{%
  \tikzexternaldisable\tikz[overlay,remember picture,baseline] \node [anchor=base] (#1) {};\tikzexternalenable}

\newcommand\MyLine[3][]{%
  \tikzexternaldisable\begin{tikzpicture}[overlay,remember picture]
    \draw[#1] (#2.south) -- (#3.north);
  \end{tikzpicture}\tikzexternalenable}
  
%--Jordankasten
\newcommand{\Jord}{ \tikzexternaldisable \begin{BMAT}{|ccc|}{|ccc|} %
				0 & & \\ %
				1\tikzmark{z} &  \diagdown & \\ %
				& \tikzmark{u}1 \MyLine[thick]{z}{u} &  0 %
			\end{BMAT} \tikzexternalenable}
			
%--Jordankasten mit Lambda
\newcommand{\JordLa}[1][\empty]{%
	\tikzexternaldisable\ifthenelse{\equal{#1}{\empty}}
		{\begin{BMAT}{|ccc|}{|ccc|} %
				\lambda & & \\ %
				1\tikzmark{z} &  \diagdown & \\ %
				& \tikzmark{u}1 \MyLine[thick]{z}{u} &  \lambda  %
			\end{BMAT}}
		{\begin{BMAT}{|ccc|}{|ccc|} %
				\lambda_{#1}  & & \\ %
				1\tikzmark{z} &  \diagdown & \\ %
				& \tikzmark{u}1 \MyLine[thick]{z}{u} &  \lambda_{#1}  %
			\end{BMAT}}\tikzexternalenable}

%--Inhaltsverzeichnis
\usepackage[tocindentauto]{tocstyle}
\usetocstyle{KOMAlike}			

\newcommand{\fach}{Einführung in die Funktionalanalysis}
\newcommand{\semester}{WiSe 2014}
\newcommand{\homepage}{https://wwwmath.uni-muenster.de/u/wilhelm.winter/wwinter/funktionalanalysis.html}
\newcommand{\prof}{Prof.\,Dr.\,Wilhelm Winter}
\providecommand{\verfasser}{Jannes Bantje}

%--Konfiguration von scrheadings
\setheadsepline{1pt}[\color{light_gray}]
\pagestyle{scrheadings}
%\fancyhf[H, F]{}
\clearscrheadfoot

\providecommand{\shortFach}{\fach}
\lehead{ \includegraphics[height=0.5 cm,keepaspectratio]{../!config/Bilder/Logo_WWU_Muenster_light_gray.pdf}}
\rehead{\rule{0cm}{0.5cm}\footnotesize \sffamily \color{light_gray} \verfasser -- Mitschrift \shortFach}
\lohead{\rule{0cm}{0.5cm} \footnotesize \sffamily \color{light_gray} Stand: \today \; \thistime[:]}
\rohead{\includegraphics[height=0.5 cm,keepaspectratio]{../!config/Bilder/fb10logo_gray.pdf}}


\ofoot[{ \color{dark_gray} \LARGE \sffamily \thepage}]{{ \color{dark_gray} \LARGE \sffamily \thepage}} %hier wir auch der plain Stil bearbeitet!
\automark{section}
\ifoot{ \color{dark_gray} \small \leftmark}

%--Metadaten
\author{\verfasser}
\titlehead{\includegraphics[height=1.5cm, keepaspectratio]{../!config/Bilder/Logo_WWU_Muenster.pdf}%
\hfill \includegraphics[height=1.3cm, keepaspectratio]{../!config/Bilder/fb10logo.pdf}}
\title{Skript \fach}
\subtitle{Mitschrift der Vorlesung  \enquote{\fach} von \prof}

\publishers{
Aktuelle Version verfügbar bei: \bigskip\\
\normalsize
\begin{minipage}[t]{0.45\textwidth}
	\centering{
	\qrcode[height=3.5cm, version=4]{https://github.com/JaMeZ-B/latex-wwu}
	\medskip\\
	\includegraphics[height=0.4cm, keepaspectratio]{../!config/Bilder/github_octo.pdf} 
	\includegraphics[height=0.4cm, keepaspectratio]{../!config/Bilder/GitHub_Logo.pdf} (inklusive Sourcecode)\\
	{\small\url{https://github.com/JaMeZ-B/latex-wwu}}}
\end{minipage}
\hfill
\begin{minipage}[t]{0.45\textwidth}
	\centering{
	\qrcode[height=3.5cm, version=4]{btsync://B6WH2DISQ5QVYIRYIEZSF4ZR2IDVKPN3I?n=latex_share}
	\medskip\\
	\includegraphics[height=0.4cm, keepaspectratio]{../!config/Bilder/bt_sync_logo.pdf} 
	{\large \textbf{Bittorrent} Sync} \\
	\texttt{B6WH2DISQ5QVYIRYIEZSF4ZR2IDVKPN3I}
	}
\end{minipage}
}



\begin{document}
\pagenumbering{Roman}
\maketitle
\begin{abstract}
\section*{Vorwort --- Mitarbeit am Skript}
Dieses Dokument ist eine Mitschrift aus der Vorlesung \enquote{\fach, \semester}, gelesen von \prof. Der Inhalt entspricht weitestgehend dem Tafelanschrieb. Für die
Korrektheit des Inhalts übernehme ich keinerlei Garantie! Für Bemerkungen und Korrekturen -- und seien es nur Rechtschreibfehler -- bin ich sehr dankbar. 
Korrekturen lassen sich prinzipiell auf drei Wegen einreichen: 
\begin{itemize}
	\item Persönliches Ansprechen in der Uni, Mails an \hrefsymmail{mailto:\mail}{\mail} (gerne auch mit annotieren PDFs) 
	\item \emph{Direktes} Mitarbeiten am Skript: Den Quellcode poste ich auf GitHub (siehe oben), also stehen vielfältige Möglichkeiten der Zusammenarbeit zur Verfügung:
	Zum Beispiel durch Kommentare am Code über die Website und die Kombination Fork + Pull Request. Wer sich verdient macht oder ein Skript zu einer Vorlesung, die 
	ich nicht besuche, beisteuern will, dem gewähre ich gerne auch Schreibzugriff.
	
	Beachten sollte man dabei, dass dazu ein Account bei \url{github.com} notwendig ist, der allerdings ohne Angabe von persönlichen Daten angelegt werden kann. 
	Wer bei GitHub (bzw. dem zugrunde liegenden Open-Source-Programm \enquote{\texttt{git}}) -- verständlicherweise -- Hilfe beim Einstieg braucht, dem helfe ich gerne 
	weiter. Es gibt aber auch zahlreiche empfehlenswerte Tutorials im Internet.\footnote{zB. \url{https://try.github.io/levels/1/challenges/1}, ist auf Englisch, aber dafür 
	interaktives LearningByDoing}
	\item \emph{Indirektes} Mitarbeiten: \TeX-Dateien per Mail verschicken. 
	
	Dies ist nur dann sinnvoll, wenn man einen ganzen Abschnitt ändern möchte (zB. einen alternativen Beweis geben), da ich die Änderungen dann per Hand einbauen muss! Ich freue mich aber auch über solche Beiträge!
\end{itemize}
\section*{Vorlesungshomepage}
{\centering 
\begin{minipage}[c][][c]{\textwidth}
	\centering \qrcode[height=3cm]{\homepage} \medskip\\
	\footnotesize \url{\homepage}
\end{minipage}
\par}

\end{abstract}

\setcounter{page}{1}
\tableofcontents
\cleardoubleoddemptypage
\pagenumbering{arabic}
\setcounter{page}{1}

\section{Metrische Räume und der Satz von Baire} % (fold)
\label{sec:1}

\subsection[Definition: Metrischer Raum]{Definition} % (fold)
\label{sub:11}
Ein \Index{metrischer Raum} ist ein Paar $(X,d)$, wobei $X$ eine Menge und $d : X \times X \to [0,\infty)$ ist, sodass\marginnote{14 Okt}
\begin{enumerate}[1)]
	\item $d(x,y)=0 \iff x=y$ $\forall x,y \in X$
	\item $d(x,y) = d(y,x)$ $\forall x,y \in X$
	\item $d(x,z) \le d(x,y) + d(y,z)$ $\forall x,y,z \in X$
\end{enumerate}
% subsection 11 (end)

\subsection[Definition: Offen, abgeschlossen und Abschluss]{Definition} % (fold)
\label{sub:12}
Sei $(X,d)$ ein metrischer Raum
\begin{itemize}
	\item Eine Teilmenge $U \subseteq X$ heißt \Index{offen}, falls für jedes $x \in U$ ein $\varepsilon > 0$ existiert, so dass 
	\[
		B(x,\varepsilon) := \set[y \in X]{d(x,y) < \varepsilon} \subset U 
	\]
	\item Eine Teilmenge $A \subset X$ heißt \Index{abgeschlossen}, falls $X \setminus A$ offen ist (als Teilmenge von $X$). 
	\item $\mathcal{T}X := \set[U \subset X]{U \text{ offen}} $ ist die Topologie auf $X$ (die von der Metrik $d$ induziert wird)
	\item Falls $W \subset X$ eine Teilmenge ist, dann bezeichnet $\overline{W}$ den \Index{Abschluss} von $W$, d.h. die kleinste abgeschlossene Teilmenge von $X$, die 
	$W$ enthält. Es gilt
	\[
		\overline{W} = \enspace\bigcap_{\mathclap{A \subset X \text{abg.}, W \subset A}} \enspace A 
	\] 
\end{itemize}
Für metrische Räume gilt:
\(
	\overline{W} = \set[\lim_n x_n]{(x_n)_n \subset W \text{ konvergente Folge}}.
\)
Warum gilt $W \subset \overline{W}$?
% subsection 12 (end)

\subsection[Definition: Stetigkeit, gleichmäßige Stetigkeit, Isometrie]{Definition} % (fold)
\label{sub:13}
Eine Abbildung $f : X \to Y$ zwischen zwei metrischen Räumen $(X,d_X)$ und $(Y,d_Y)$ heißt 
\begin{itemize}
	\item \bet{stetig in} $x \in X$, falls
	\(
		\forall \varepsilon >0 : \exists \delta >0 : \forall x' \in X : d_X(x,x') < \delta  \Longrightarrow d_Y \enbrace[\big]{f(x), f(x')} < \varepsilon 
	\)
	\item \Index{stetig}, falls $f$ an jedem Punkt $x \in X$ stetig ist.
	
	(Äquivalent: Für jede offene Menge $V$ in $Y$ ist $f ^{-1}(V)$ offen in $X$) 
	\item \Index{gleichmäßig stetig}, falls
	\(
		\forall \varepsilon>0 : \exists \delta >0 : \forall x,x' \in X : d_X(x,x') < \delta \Longrightarrow d_Y \enbrace[\big]{f(x), f(x')} < \varepsilon .
	\) 
	\item $f$ heißt \Index{Isometrie}, falls \marginnote{automatisch injektiv}
	\(
		\forall x,x' \in X : d_X(x,x') = d_Y \enbrace[\big]{f(x), f(x')}.
	\)
	\item $f$ heißt \Index{isometrischer Isomorphismus}, falls $f$ bijektiv und isometrisch ist. \marginnote{$f ^{-1}$ auch}
\end{itemize}
% subsection 13 (end)

\subsection[Definition: Cauchy-Folge und Vollständigkeit]{Definition} % (fold)
\label{sub:14}
Eine Folge $(x_n)_{n \in \mathds{N}}$ in einem metrischen Raum $(X,d)$ heißt \bet{Cauchy}\index{Cauchy-Folge}, falls
\[
	\forall \varepsilon>0 : \exists N \in \mathds{N} : \forall n,k >N : d(x_n,x_k) < \varepsilon.
\]
$(X,d)$ heißt \Index{vollständig}, falls jede Cauchy-Folge in $X$ konvergiert. 
% subsection 14 (end)

\subsection[Satz: Existenz einer eindeutigen Vervollständigung metrischer Räume]{Satz} % (fold)
\label{sub:15}
Sei $(X,d)$ ein metrischer Raum. Dann existiert ein vollständiger metrischer Raum $(\tilde{X}, \tilde{d})$ und eine Isometrie $\iota : X \hookrightarrow \tilde{X}$, sodass
$\overline{\iota(X)} = \tilde{X}$ (d.h. $\iota(X)$ ist dicht in $\tilde{X}$). $(\tilde{X}, \tilde{d})$ heißt \Index{Vervollständigung} von $(X,d)$ und ist eindeutig bis auf
isometrische Isomorphie.
\minisec{Beweis}
\begin{description}
	\item[Eindeutigkeit:] 
	Angenommen, $(\hat{X}, \hat{d})$ ist ein weiterer vollständiger metrischer Raum und $\kappa : X \to \hat{X}$ eine Isometrie mit 
	$\overline{\kappa(X)}=\hat{X}$. Definiere $\gamma : \hat{X} \to \tilde{X}$ wie folgt: Falls $y \in \hat{X}$, wähle eine Folge $(x_n)_{n \in \mathds{N}}$ in $X$, sodass
	$y= \lim_{n \to \infty} \kappa(x_n)$. Setze nun
	\[
		\gamma(y) := \lim_{ n \to \infty} \iota(x_n) \in \tilde{X} 
	\]
	Zu zeigen: $\gamma$ ist ein isometrischer Isomorphismus.\todo{Achtung Wortspiel: Beweis \emph{vervollständigen} ;-D}
	\begin{description}
		\item[Injektivität:] Seien $y,y' \in \hat{X}$ mit $\gamma(y)=\gamma(y')$. Dann existieren Folgen $(x_n)_n, (x_n')_n$ in $X$ mit
		\[
			\lim_{ n \to \infty} \iota(x_n) = \lim_{ n \to \infty} \iota(x_n') 
		\]
	\end{description}
	\item[Existenz:] \begin{description}
		\item[Konstruktion von $(\tilde{X}, \tilde{d})$:] Setze $Y := \set[(x_n)_{n \in \mathds{N}}]{(x_n)_n \text{ ist Cauchy-Folge in }X}$. Definiere 
	\[
		(x_n)_n \sim (x_n') :\iff \lim_{ n \to \infty} d(x_n, x_n') = 0
	\]
	$\sim$ ist eine Äquivalenzrelation auf $Y$. Definiere nun $\tilde{X} := \nicefrac{Y}{\sim}$ und $\tilde{d} :  \tilde{X} \times \tilde{X} \to [0, \infty)$ durch
	\[
		\tilde{d} \enbrace*{\benbrace[\big]{(x_n)_n} , \benbrace[\big]{(x_n')_n}  } := \lim_{ n \to \infty} d(x_n , x_n') 
	\]
	$\tilde{d}$ ist eine wohldefiniert, d.h. falls $\benbrace*{(x_n)_n} = \benbrace*{(y_n)_n} $ und $\benbrace*{(x_n')_n} = \benbrace*{(y_n')_n}$, dann ist
	\[
		\lim_{n \to \infty} d(x_n , x_n') = \lim_{n \to \infty} d(y_n, y_n'). \tag*{(leichte Übung)}
	\]
	Weiter ist $\enbrace[\big]{d(x_n x_n')}_n$ eine Cauchy-Folge in dem vollständigen Raum $[0,\infty)$ und somit 
	konvergent: Sei dazu $\varepsilon>0$. Dann existiert ein $N \in \mathds{N}$, sodass $d(x_n,x_m) \le \frac{\varepsilon}{2}$ und $d(x_n', x_m') \le \frac{\varepsilon}{2}$ für alle $n,m \ge N$. Dann gilt nach der \hyperref[sub:vier_ungl]{Vierecksungleichung (siehe Anhang \ref*{sub:vier_ungl})}
	\[
		\abs[\big]{d(x_n,x_n') - d(x_m,x_m')} \le d(x_n, x_m) + d(x_n', x_m') \le \varepsilon
	\]
	\item[Einbettung von $X:$] Definiere nun $\iota : X \to \tilde{X}$ durch $x \mapsto [(x,x,x, \ldots )] \in \tilde{X}$. $\iota$ ist Isometrie, da
	\[
		\tilde{d} \enbrace[\big]{\iota(x),\iota(y)} = \lim_{ n \to \infty} d(x,y) = d(x,y)
	\]
	% Falls $\benbrace*{(x_n)_n} \in \tilde{X}$, finde $y_k \in \iota(X)$, so dass $\benbrace*{(x_n)} = \lim_{ k} y_k$ und $y_k = \benbrace*{(x_k, x_k, \ldots )}= \iota(x_k)$.
% 	Dann gilt $y_k \xrightarrow{k \to \infty} \benbrace*{(x_n)_n} $, d.h.
% 	\[
% 		\lim_{ n \to \infty} d(x_k,x_n) = \tilde{d} \enbrace*{y_k, \benbrace*{(x_n)_n}} \xrightarrow{k \to \infty} 0
% 	\]
% 	$\leadsto \tilde{X}= \overline{\iota(X)}$.
	Sei nun $\benbrace*{(x_n)_n} \in \tilde{X}$ und $\varepsilon>0$. Da $(x_n)_n$ eine Cauchy-Folge ist, gibt es ein $N \in \mathds{N}$, sodass für alle
	$n,m \ge N$ gilt $d(x_n,x_m) <\varepsilon$. Dann gilt
	\[
		\tilde{d}\enbrace[\big]{\iota(x_N), \benbrace*{(x_n)_n} } = \lim_{ n \to \infty} d \enbrace*{x_N, x_n} < \varepsilon
	\]
	\item[Vollständigkeit von $(\tilde{X}, \tilde{d})$:] Sei $\enbrace*{\overline{x}^m}_m$ eine Cauchyfolge in $\tilde{X}$.\marginnote{$\overline{x} \in \tilde{X}$}
	\[
		\Longrightarrow \forall \varepsilon >0 : \exists M(\varepsilon) \in \mathds{N} : \forall m,m' > M : \tilde{d} \enbrace*{\overline{x}^m, \overline{x}^{m'}} < \frac{\varepsilon}{3} 
	\]
	Wenn $\overline{x}^m = \benbrace*{(x^m_n)_n}$ und $\overline{x}^{m'} = \benbrace*{(x^{m'}_n)_n} $, dann gilt also für alle $m,m' >M(\varepsilon)$
	\begin{equation*}
		\lim_{ n \to \infty} d \enbrace*{x^m_n, x^{m'}_n} < \frac{\varepsilon}{3}  \label{eq:15:absch1}\tag{*}
	\end{equation*}
	Für alle $m$ ist $(x^m_n)_n$ eine Cauchyfolge, also gilt
	\begin{equation*}
		\forall m : \exists N(m) : \forall n,n' \ge N(m) : d \enbrace*{x^m_n, x^m_{n'}} < \frac{1}{m}  \label{eq:15:absch2}\tag{**}
	\end{equation*}
	Setze nun $z_n := x^n_{N(n)}$. Behauptung: $(z_n)_n$ ist eine Cauchyfolge. Sei $\varepsilon>0$. Dann gilt für $n,m>N=\max\set{M(\varepsilon), \frac{3}{\varepsilon}}$
	für ein $k > N(m), N(n)$ \marginnote{$d(z_n,z_m)$ ist unabhängig von $k$, also kann man $k$ beliebig groß wählen}
	\begin{align*}
		d(z_n,z_m) = d \enbrace*{x^n_{N(n)}, x^m_{N(m)}} \le 
		\underbrace{d\enbrace*{x^n_{N(n)}, x^n_k}}_{\stackrel{\text{\eqref{eq:15:absch2}}}{<} \frac{1}{n} < \frac{1}{N} \le \frac{\varepsilon}{3}} + 
		\underbrace{d\enbrace[\Big]{x^n_k, x^m_k}}_{\stackrel{\text{\eqref{eq:15:absch1}}}{\le} \frac{\varepsilon}{3}}   + 
		\underbrace{d\enbrace*{x^m_k, x^m_{N(m)}}}_{\stackrel{\text{\eqref{eq:15:absch2}}}{<} \frac{1}{m} < \frac{1}{N} \le \frac{\varepsilon}{3}} 
		< \varepsilon
	\end{align*}
	$\Rightarrow (z_n)_n$ ist eine Cauchyfolge, also 
	\begin{equation*}
		\forall \varepsilon>0 : \exists N_z(\varepsilon) : \forall n,m > N_z(\varepsilon)  : d(z_n,z_m) < \varepsilon  \label{eq:15:absch3} \tag{\#}
	\end{equation*}
	Es bleibt zu zeigen: $\lim_{m \to \infty} \overline{x}^m = \benbrace*{(z_n)_n}$. Sei dazu $\varepsilon>0$. Dann gilt für 
	$m \ge \max \set{\frac{2}{\varepsilon}, N_z(\varepsilon)}$
	\begin{align*}
		\tilde{d} \enbrace[\big]{\overline{x}^m, \benbrace*{(z_n)_n}} = \lim_{ n \to \infty} d \enbrace*{x^m_n, x^n_{N(n)}}\le \lim_{ n \to \infty} 
		\enbrace[\Bigg]{ \underbrace{d\enbrace*{x_n^m, x^m_{N(m)}}}_{\stackrel{\text{\eqref{eq:15:absch2}}}{<} \frac{1}{m} \le \frac{\varepsilon}{2}  } +
		\underbrace{d\enbrace*{x^m_{N(m)}, x^n_{N(n)}}}_{=d(z_m,z_n) \stackrel{\text{\eqref{eq:15:absch3}}}{<} \frac{\varepsilon}{2} } } < \varepsilon
	\end{align*}
	Also gilt $\overline{x}^m \xrightarrow{m \to \infty} \benbrace*{(z_n)_n}$ und $(\tilde{X},\tilde{d})$ ist vollständig. \bewende
	\end{description}
\end{description}
% subsection 15 (end)

\subsection[Definition: Raum der beschränkten, stetigen Abbildungen]{Definition} % (fold)
\label{sub:16}
Sei $(W,\mathcal{T})$ ein topologischer Raum und $(X,d)$ ein metrischer Raum. Sei 
\[
	C_b(W,X) = \set[f : W \to X]{f \text{ stetig und beschränkt}} 
\]
versehen mit der Metrik $d_{W,X}$, definiert durch
\[
	d_{W,X}(f,g) = \sup_{t \in W} d \enbrace[\big]{f(t), g(t)} 
\]
% subsection 16 (end)

\subsection[Bemerkung: $d_{W,X}$ als Metrik auf $C(W,X)$]{Bemerkung} % (fold)
\label{sub:17}
Auf $C(W,X) = \set{f : W \to X \text{ stetig}} $ ist $d_{W,X}$ eine \Index{erweiterte Metrik}, d.h. der Wert $\infty$ ist möglich.
$\tilde{d}_{W,X} := \min \set{1, d_{W,X}} $ ist eine \enquote{echte} Metrik auf $C(W,X)$.
% subsection 17 (end)

\subsection[Proposition: $X$ vollständig $\Rightarrow$ $C-B(W,X)$ und $C8W,X$ vollständig]{Proposition} % (fold)
\label{sub:18}
Falls $X$ vollständig ist, dann sind $C_b(W,X)$ und $C(W,X)$ vollständig (bezüglich $d_{W,X}$ bzw. $\tilde{d}_{W,X}$).
\minisec{Beweis}
Sei $(f_n)_n$ eine Cauchy-Folge in $C_b(W,X)$, also
\[
	\forall \varepsilon>0 : \exists N \in \mathds{N} : \forall n,m >N : \sup_{t \in W} d \enbrace[\big]{f_n(t), f_m(t)}< \varepsilon 
\]
$\Rightarrow$ für alle $t$ ist $\enbrace[\big]{f_n(t)}_n$ eine Cauchyfolge in $X$. Da $X$ vollständig ist, existiert ein $x_t \in X$ sodass 
$f_n(t) \xrightarrow{n \to \infty} x_t$. Definiere $f(t) := x_t$ punktweise.
\begin{description}
	\item[Vorbereitung:] Es gilt $\forall \varepsilon>0 : \exists N_t$ sodass $d(x_t, f_{N_t}(t)) \le \varepsilon$. Weiter ist
	\[
		\sup_{t \in W} d \enbrace[\big]{f_n(t), f_m(t)}\le \varepsilon \enspace \forall n,m \ge N_t
	\]
	$\Rightarrow \forall t : d \enbrace*{f_n(t), f_m(t)}\le \varepsilon \enspace \forall n,m \ge N$
\end{description}
\todo{Beweis zu Ende führen}
% subsection 18 (end)

\subsection[Proposition über eine Folge von Bällen]{Proposition} % (fold)
\label{sub:19}
Sei $(X,d)$ ein vollständiger metrischer Raum und $\enbrace*{\overline{B}(x_n, \varepsilon_n)}_{n \in \mathds{N}}$ mit 
$\overline{B}(x_{n+1}, \varepsilon_{n+1}) \subseteq \overline{B}(x_n,\varepsilon_n)$ und $\varepsilon_n \xrightarrow{n \to \infty} 0$. Dann existiert genau 
ein Punkt in $\bigcap_{n \in \mathds{N}} \overline{B}(x_n, \varepsilon_n)$.
\minisec{Beweis}
\begin{description}
	\item[Existenz:] Behauptung: Die Folge der Mittelpunkte $(x_n)_{n \in\mathds{N}}$ ist eine Cauchy-Folge: Sei $\varepsilon >0$. Finde $N \in \mathds{N}$,
	sodass $\varepsilon_n \le \varepsilon$ für alle $n \ge N$. Dann gilt für alle $n \ge m \ge N$
	\[
		d(x_n, x_m) \le \varepsilon_m \le \varepsilon,
	\]
	da $x_n \in \overline{B}(x_m, \varepsilon)$ ist. Da $X$ vollständig ist, existiert ein $x \in X$ mit $x_n \xrightarrow{n \to \infty} x$.
	
	Behauptung: $x \in \bigcap_{n \in \mathds{N}} \overline{B}(x_n, \varepsilon_n)$. Wähle dazu ein $N \in \mathds{N}$. Dann ist $(x_n)_{n \ge N}$ eine
	Cauchy-Folge in $\overline{B}(x_N,\varepsilon_N)$. $\overline{B}(x_N,\varepsilon_N)$ ist ein abgeschlossener Teilraum von $X$ und somit vollständig. 
	Also ist $x \in \overline{B}(x_N,\varepsilon_N)$. Da $N$ beliebig war, gilt $x \in \overline{B}(x_n,\varepsilon_n)$ für jedes $n \in \mathds{N}$.
	\item[Eindeutigkeit:] Es seien $x,y \in \bigcap_{n \in \mathds{N}} \overline{B}(x_n,\varepsilon_n)$, also 
	$x,y \in \overline{B}(x_n, \varepsilon_n) $ für alle $ n \in \mathds{N}$. Dann gilt
	\[
		d(x,y) \le d(x,x_n) + d(x_n,y) \le 2 \cdot \varepsilon_n \xrightarrow{n \to \infty} 0
	\]
	für jedes $n \in \mathds{N}$ und somit $d(x,y)=0 \Leftrightarrow x=y$. \bewende
\end{description}
% subsection 19 (end)

\subsection[Satz von Baire]{Satz von Baire\protect\footnote{nach René Louis Baire, \url{https://de.wikipedia.org/wiki/René_Louis_Baire}}} % (fold)
\label{sub:110} 
Es gelten folgende äquivalente Formulierungen:
\begin{enumerate}[a)]
	\item Sei $(X,d)$ ein vollständiger metrischer Raum und $A_0, A_1, \ldots $ eine Folge abgeschlossener Teilmengen. Falls $\bigcup_{n \in \mathds{N}} A_n$ eine
	offene Kugel enthält, so auch eines der $A_n$.
	\item In einem vollständigen metrischen Raum hat eine abzählbare Vereinigung von abgeschlossenen Mengen ohne innere Punkte keine inneren Punkte.
	\item In einem vollständigen metrischen Raum ist ein abzählbarer Durchschnitt von dichten offenen Mengen wieder dicht.
\end{enumerate}
\minisec{Beweis}
Zunächst die Äquivalenz: 
\begin{description}
	\item[a)$\Rightarrow$b):] Klar, da b) Kontraposition von a) ist.
	\item[b)$\Rightarrow$c):] Sei $U_n \subseteq X$ offen und dicht $\Rightarrow X \setminus U_n \subseteq X$ abgeschlossen und hat somit keine inneren Punkte. Aus b) folgt
	nun 
	\[
		\bigcup_{n \in \mathds{N}}X \setminus U_n = X \setminus \bigcap_{n \in \mathds{N}} U_n
	\]
	hat keine inneren Punkte. Also ist $\bigcap_{n \in \mathds{N}} U_n$ dicht in $X$.
	\item[c)$\Rightarrow$a):] ebenso.
\end{description}
Wir wollen nun a) durch Widerspruch beweisen, d.h. wir nehmen an, dass gilt: Jede offene Kugel schneidet $X \setminus A_k$ für jedes $k \in \mathds{N}$. Dazu wollen wir 
Folgen $(x_k)_{k \in \mathds{N}} \subset X$, $(\varepsilon_k)_{k \in \mathds{N}} \subset (0,1]$ finden mit 
\begin{enumerate}[(i)]
	\item $\varepsilon_k < \frac{1}{k+1}$,
	\item $\overline{B}(x_{k+1}, \varepsilon_{k+1}) \subset (X \setminus A_k) \cap B(x_k, \varepsilon_k) \subset \overline{B}(x_k, \varepsilon_k)$,
	\item $\overline{B}(x_k, \varepsilon_k) \subset \bigcup_{n \in \mathds{N}} A_n$
\end{enumerate}
für $k \in \mathds{N}$. Dann gilt 
\[
	\bigcup_{n \in \mathds{N}} A_n \stackrel{\text{(iii),(ii)}}{\supset} \bigcap \overline{B}(x_k, \varepsilon_k) \stackrel{\text{(ii)}}{\subset} 
	\bigcap_{k \in \mathds{N}}(X \setminus A_k) \cap B(x_k, \varepsilon_k) \subset \bigcap_{k \in \mathds{N}} X \setminus A_k = X \setminus 
	\enbrace*{\bigcup_{k \in \mathds{N}} A_k} 
\]
Aber wegen (i) und (ii) existiert nach Proposition \ref{sub:19} ein $x \in \bigcap_{k \in \mathds{N}}\overline{B}(x_k,\varepsilon_k)$. \light \medskip \\
Wir suchen also eine Abbildung $\overline{c} : \mathds{N} \to X \times (0,1], k  \mapsto(x_k, \varepsilon_k)$ mit (i),(ii),(iii) für $k \in \mathds{N}$. Setze
\begin{align*}
	P_m &:= \set[{c : \set{0, \ldots ,m} \to X \times (0,1] }]{\begin{array}{l}
		\text{(i),(iii) gilt für }k \in \set{0, \ldots , m}, \\
		\text{(ii) gilt für }k \in \set{0, \ldots , m-1}
	\end{array} }  \\
	P_\infty &:= \set[{c : \mathds{N} \to X \times (0,1]}]{\text{(i),(ii),(iii) gilt für } k \in \mathds{N} \rule{0cm}{0.3cm}} 
\end{align*}
Die Menge $P := \enbrace*{\bigcup_{m \in \mathds{N}} P_m} \cup P_\infty$ ist partiell geordnet bezüglich $\prec$:
\[
	c \prec c' \text{ falls } m \le m' \text{ und } c'\big|_{\set{0,\ldots ,m}} =c , \text{ bzw.} c'=c \text{ falls } m=m'=\infty
\]
$P$ ist nicht leer, denn nach Vorraussetzung existiert $(x_0,\varepsilon_0)$ mit $0 < \varepsilon_0 < 1$ und $\overline{B}(x_0, \varepsilon_0) \subset 
\bigcup_{n \in\mathds{N}} A_n$, d.h. $0 \mapsto (x_0,\varepsilon_0) \in P_0$. Jede total geordnete Teilmenge $\emptyset \not= \Gamma$ von $P$ besitzt eine obere Schranke:
\begin{itemize}
	\item Falls $\Gamma$ ein $c\in P_\infty$ enthält, so ist $c$ obere Schranke.\marginnote{warum?}
	\item Ebenso falls $\Gamma$ ein $c \in P_{\overline{m}}$ enthält und $\Gamma \cap P_{m'}= \emptyset$ für alle $\overline{m} < m' \in \mathds{N} \cup \set{\infty}$.
	\item Falls $\Gamma \subset \bigcup_{\mathds{N}} P_m$, aber $\Gamma \not\subset \bigcup_{m \le \overline{m}} P_m$ für jedes $\overline{m}$, so definieren wir eine
	obere Schranke in $P_\infty$ durch Einschränkung.\marginnote{wie genau?} 
\end{itemize}
Mit dem Lemma von Zorn folgt, dass $P$ ein maximales Element $\overline{c}$ besitzt. Behauptung: $\overline{c}\in P_\infty$ wie gewünscht:

Falls $\overline{c} \in P_m$ für ein $m \in \mathds{N}$, so gilt nach Annahme, dass $\enbrace*{X \setminus A_m} \cap B(x_m, \varepsilon_m) \not= \emptyset$. Dann 
existiert aber $(x_{m+1},\varepsilon_{m+1}) \in X \times (0,1]$ mit $\varepsilon_{m+1} < \frac{1}{m+2}$ und 
\[
	\overline{B}(x_{m+1}, \varepsilon_{m+1}) \subset  \enbrace*{X \setminus A_m} \cap B(x_m,\varepsilon_m) 
\]
Definiere $\overline{\overline{c}} : \set{0, \ldots , m+1} \to X \times (0,1]$ durch
\[
	k \longmapsto \begin{cases}
		(x_k,\varepsilon_k), &\text{ falls } k=m+1\\
		\overline{c}(k), &\text{ falls } k \in \set{0, \ldots , m} 
	\end{cases},
\]
dann gilt $\overline{c} \prec \overline{\overline{c}} \in P_{m+1}$. \light zur Maximalität von $\overline{c}$. \bewende
% subsection satz_baire (end)

\subsection[Bemerkungen zum Satz von Baire]{Bemerkungen} % (fold)
\label{sub:111}
\begin{enumerate}[(i)]
	\item Die Aussage gilt auch für lokalkompakte Hausdorffräume. \hfill(Übung)
	\item Tatsächlich genügt eine schwächere Form des Auswahlaxioms(DC); das abzählbare Auswahlaxiom jedoch nicht.
	\item Falls $X$ \Index{separabel} ist (d.h. falls eine abzählbare dichte Teilmenge von $X$ existiert), dann lässt sich der Satz auch ohne (AC) beweisen.\hfill (Übung)
\end{enumerate}
% subsection 111 (end)

\subsection{Korollar: Satz der gleichmäßigen Beschränktheit} % (fold)
\label{sub:112}
Sei $(X,d)$ ein vollständiger metrischer Raum. Sei $F \subset C(X,\mathds{R})$ eine Menge, die \Index{punktweise gleichmäßig beschränkt} ist, d.h. für jedes $x \in X$ 
existiert $K_x \in \mathds{R}$, sodass
\[
	\abs*{f(x)} \le K_x \enspace \forall f \in F. 
\] 
Dann existieren $\emptyset \not= U \subset X$ offen und $K \in \mathds{R}$ so, dass 
\[
	\abs*{f(x)} \le K \enspace \forall x \in U, f \in F 
\]
\minisec{Beweis}
Definiere 
\begin{align*}
	X \supset A_n := \set[x \in X]{\abs*{f(x)} \le n \forall f \in F \rule{0cm}{0.3cm}} = \bigcap_{f \in F} f ^{-1} \enbrace[\big]{[-n,n]} 
\end{align*}
Also sind die $A_n \subset X$ abgeschlossen. Es ist $\bigcup_{n \in \mathds{N}} A_n = X$, wegen der Voraussetzung von punktweiser gleichmäßiger Beschränktheit. Außerdem enthält dies eine nichtleere offene Menge. Also existiert ein $n$ und $U \subset X$ nichtleer und offen mit $U \subset A_n$. \bewende
% subsection 112 (end)
% section 1 (end)
\newpage

\section{Normierte Räume, Hahn-Banach Sätze} % (fold)
\label{sec:2}

\subsection[Definition: Topologischer Vektorraum]{Definition} % (fold)
\label{sub:21}
Sei $X$ ein $\mathds{K}$-Vektorraum mit einer Topologie $\mathcal{T}$. Wir sagen $X$ ist ein \Index{topologischer Vektorraum}, falls Addition und Skalarmultiplikation stetig sind:
\begin{align*}
	+ : X \times X  \to X \quad &, \quad(x,y) \mapsto x+y \\
	\cdot : \mathds{K} \times X \to X \quad &, \quad (\lambda ,x) \mapsto \lambda \cdot x
\end{align*}
% subsection 21 (end)

\subsection[Proposition: Unterräume topologischer Vektorräume sind topologische Vektorräume]{Proposition} % (fold)
\label{sub:22}
Sei $X$ ein topologischer $\mathds{K}$-Vektorraum und $Y \subset X$ ein Untervektorraum. Dann sind $Y$ und $\overline{Y}$ topologische Vektorräume mit der Unterraumtopologie.
\minisec{Beweis}
\begin{itemize}
	\item Klar für $Y$.
	\item Zu zeigen: $\overline{Y}$ ist ein Untervektorraum. Seien $x_0, y_0 \in \overline{Y}$ und sei $U$ eine offene Menge in $X$ mit $x_0 + y_0 \in U$. Dann existieren
	offene Mengen $V,W$ von $X$ mit $x_0 \in V$, $y_0 \in W$ und $V+W \subset U$, da die Addition stetig ist. Da $x_0 \in \overline{Y}$ und $x_0 \in V$, existiert 
	$x_1 \in Y$ mit $x_1 \in V$. Analog existiert $y_1 \in Y$ mit $y_1 \in W$. Daher ist $x_1 + y_1 \in Y \cap U$. Da $U$ eine beliebige offene Umgebung um $x_0+ y_0$ ist,
	folgt $x_0 +y_0 \in \overline{Y}$.
	
	Skalarmultiplikation genauso. \bewende
\end{itemize}
% subsection 22 (end)

\subsection[Proposition: Normierte Vektorräume sind topologische Vektorräume]{Proposition} % (fold)
\label{sub:23}
Ein normierter Vektorraum $(X, \norm{\cdot})$ ist ein topologischer Vektorraum bezüglich der von $\norm{\cdot}$ induzierten Topologie.
\minisec{Beweis}
Es gilt 
\begin{align*}
	\norm[X]{(x+y) - (x'+y')} \le \norm[X]{x-x'} + \norm[X]{y-y'}   
\end{align*}
Daher ist $+ : X \times X \to X$ gleichmäßig stetig bezüglich der Norm $\norm[X \times X]{(z,z')} := \norm[X]{z} + \norm[X]{z'}$ auf $X \times X$. Ähnlich folgt
\begin{align*}
	\norm[X]{\lambda x - \lambda' x'} = \norm[X]{\lambda (x-x') + (\lambda -\lambda')x'} \le \abs*{\lambda } \cdot \norm[X]{x-x'} + \abs*{\lambda -\lambda'} \cdot 
	\norm[X]{x'}      
\end{align*}
Daher ist $\cdot : \mathds{K} \times X \to X$ stetig, denn falls $(\lambda_i)_{i \in I}$ ein Netz in $\mathds{K}$ und $(x_i)_{i \in I}$ ein Netz in $X$ ist mit 
$\lambda_i \to \lambda $, $x_i\to x$, dann  gilt $\lambda_i x_i \to \lambda x$, also $\norm[X]{\lambda_i x_i - \lambda x} \to 0$. \bewende
% subsection 23 (end)

\subsection[Proposition: Stetigkeit einer linearen Abbildung zwischen topologischen Vektorräumen]{Proposition} % (fold)
\label{sub:24}
Seien $X,Y$ topologische Vektorräume, $T : X \to Y$ linear. Dann sind äquivalent:
\begin{enumerate}[(1)]
	\item $T$ ist stetig.
	\item $T$ ist stetig in $0$.
	\item $T$ ist stetig in einem Punkt $\overline{x}$.
\end{enumerate}
\minisec{Beweis}
(1) $\Rightarrow$ (2) $\Rightarrow $ (3) ist klar. Wir zeigen (3) $\Rightarrow$ (1):

Für $y \in X$ definiere $L_y : X \to X$, $x \mapsto y+ x$. Dann ist $L_y$ bijektiv und stetig. Da auch $L_{-y}$ stetig ist, ist $L_y$ ein Homöomorphismus. Sei nun $T$ in 
$\overline{x}$ stetig. Sei $x_0 \in X$. Wir zeigen: $T$ ist stetig in $x_0$. Setze $y := \overline{x} - x_0$. Dann ist
\begin{align*}
	T(x_0) = T(x_0+y-y) = T(-y) + T(x_0+y) = L_{T(-y)} \circ T \circ L_y (x_0)
\end{align*}
Da $L_y(x_0) = \overline{x}$, $T$ stetig in $\overline{x}$ ist und $L_y$, $L_{T(-y)}$ stetig sind, ist somit $T$ stetig in $x_0$. \bewende
% subsection 24 (end)

\subsection[Definition: Stetige Funktionale und Operatoren]{Definition} % (fold)
\label{sub:25}
Für topologische Vektorräume $X,Y$ definieren wir 
\[
	\mathcal{L}(X,Y) = \set[T : X \to Y]{T \text{ linear und stetig}}  
\]
Wir schreiben $X':= \mathcal{L}(X,\mathds{K})$ für den \Index{Dualraum} von $X$. Die Elemente von $X'$ heißen (stetige) \Index{Funktionale}.\footnote{Aus der linearen Algebra: $X^* = \Hom(X,\mathds{R}) = \set[T : X \to Y]{T \text{ linear}}$. Die Elemente von $X^*$ heißen Funktionale}
Die Elemente von $\mathcal{L}(X,X)$ heißen (stetige) \Index{Operatoren} auf $X$. 
% subsection 25 (end)

\subsection[Bemerkung: Vektorraumstruktur auf $\mathcal{L}(X,Y)$, Algebrastruktur auf $\mathcal{L}(X,X)$]{Bemerkung} % (fold)
\label{sub:26}
$\mathcal{L}(X,Y)$ und $\mathcal{L}(X,\mathds{K}) = X'$ sind Vektorräume, wobei die Vektorraumstruktur punktweise definiert ist. $\mathcal{L}(X,X)$ ist sogar eine Algebra mittels Komposition.
% subsection 26 (end)

\subsection[Proposition: Stetigkeit linearer Abbildungen mittels Normabschätzung]{Proposition} % (fold)
\label{sub:27}
Seien $X,Y$ normierte Vektorräume, $T : X \to Y$ linear. Dann ist $T$ stetig genau dann, wenn ein $\mu \ge 0$ existiert, sodass 
\[
	\norm[Y]{T(x)} \le \mu \cdot \norm[X]{x} \quad \forall x \in X  
\]
Ein stetiger Operator zwischen normierten Vektorräumen heißt deswegen auch \bet{beschränkt}\index{Operator!beschränkter Operator}.
\minisec{Beweis}
\begin{description}
	\item[\enquote{$\Leftarrow$}:] Klar: Wenn $x_n \to 0$, dann $T(x_n) \to 0$, denn $\norm{T(x_n)} \le \mu \norm{x_n } \to 0$. Also ist $T$ stetig in $0$ und nach
	\ref{sub:24} überall.
	\item[\enquote{$\Rightarrow $}:] Angenommen $T$ ist stetig. Dann setzen wir
	\[
		\mu := \sup \set[{\frac{1}{\norm[X]{x}} \cdot \norm[Y]{T(x)}  }]{x \in X \setminus \set{0} } 
	\] 
	Falls $\mu = \infty$, dann existieren $x_n \in X$ mit $\frac{1}{\norm[X]{x_n}} \cdot \norm[Y]{T(x_n)} \ge n$. Betrachte 
	$x_n' :=  \frac{1}{n \cdot \norm[X]{x_n}} \cdot x_n $, dann $\norm{x_n'} = \frac{1}{n}$, also $x_n' \to 0$. Aber es gilt $T(x_n')\not\longrightarrow 0$, denn
	\begin{align*}
		\norm[Y]{T(x_n')} = \frac{1}{n \cdot \norm[X]{x_n} } \cdot \norm[Y]{T(x_n)} \ge 1 
	\end{align*}
	für alle $n$, \light zur Stetigkeit von $T$. Also $\mu < \infty$. \bewende \todo{Rev Chap 2}
\end{description}
% subsection 27 (end)

\subsection[Definition: Operatornorm]{Definition} % (fold)
\label{sub:28}
Seien $X,Y$ normierte Vektorräume, $T \in \mathcal{L}(X,Y)$. Die \bet{Norm} (oder \Index{Operatornorm}) von $T$ ist
\[
	\norm{T} = \sup \set[{\frac{1}{\norm[X]{x}} \cdot \norm[Y]{T(x)}}]{x \in X \setminus \set{0} } = \sup \set[{\norm[Y]{T(x)} }]{x \in X \text{ mit } \norm[X]{x} 
	= 1\rule{0cm}{0.4cm}}   
\] 
Nach \ref{sub:27} ist $\norm{T}<  \infty $ und $\norm[Y]{T x} \le \norm{T} \cdot \norm[X]{x}$ für alle $x \in X$.
% subsection 28 (end)

\subsection[Propostion: $\mathcal{L}(X,Y)$ ist ein normierter Raum mit der Operatornorm]{Proposition} % (fold)
\label{sub:29}
Seien $X,Y$ normierte Räume, betrachte $\mathcal{L}(X,Y)$ mit $\norm{.}$. Dann ist $\mathcal{L}(X,Y)$ ein normierter Raum.
\minisec{Beweis}
\begin{itemize}
	\item $\norm{T} \ge 0 $ für alle $T \in \mathcal{L}(X,Y)$.
	\item Es gilt
	\begin{align*}
		\norm{T}=0 \iff \forall x \in X, \norm[X]{x}=1 : \norm[Y]{Tx}=0 &\iff \forall x \in X, \norm[X]{x}=1 : Tx=0 \\ &\iff \forall x \in X : Tx=0 \\ &\iff T=0
	\end{align*}
	\item Sei $\lambda  \in \mathds{K}$, $T \in \mathcal{L}(X,Y)$. Dann gilt
	\[
		\norm{\lambda \cdot T} = \sup_{\norm[X]{x}=1} \norm[Y]{\lambda  \cdot T(x)}  = \sup_{\norm[X]{x}=1 } \abs*{\lambda } \cdot \norm[Y]{Tx} =
		\abs*{\lambda } \cdot \sup_{\norm[X]{x}=1 } \norm[Y]{Tx} = \abs*{\lambda } \cdot \norm{T}      
	\]
	\item Seien $T_1, T_2 \in \mathcal{L}(X,Y)$ und $x \in X$ mit $\norm[X]{x}=1 $. Dann gilt
	\[
		\norm[Y]{(T_1 + T_2)(x)} = \norm[Y]{T_1 x + T_2 x} \le \norm[Y]{T_1 x} + \norm[Y]{T_2 x}
	\]
	Damit folgt weiter
	\begin{align*}
		\norm{T_1 +T_2} = \sup_{\norm[X]{x}=1 } \norm[Y]{(T_1+ T_2) (x)} &\le \sup_{\norm[X]{x}=1} \enbrace*{\norm[Y]{T_1 x} + \norm[Y]{T_2 x}  }   \\
		&\le \sup_{\norm[X]{x}=1} \sup_{\norm[X]{x'}=1} \enbrace*{\norm[Y]{T_1 x} + \norm[Y]{T_2 x}  }  \\ &= \norm{T_1} + \norm{T_2} \bewende   
	\end{align*}
\end{itemize}
% subsection 29 (end)

\subsection[Definition: Normierte $\mathds{K}$-Algebra]{Definition} % (fold)
\label{sub:210}
Eine $\mathds{K}$-Algebra heißt \bet{normiert}\index{normierte $\mathds{K}$-Algebra}, falls $A$ mit einer Norm $\norm{.}$ versehen ist, so dass $(A,\norm{.} )$ ein normierter Vektorraum und
\[
	\norm{a b} \le \norm{a} \cdot \norm{b} \quad \forall a,b \in A   
\]
Falls $A$ unital ist, d.h. es existiert ein Einselement $1_A$, dann gilt $\norm{1_A} \le 1$.
\minisec{Bemerkung}
Wenn $\norm{1_A} < 1$, dann $\norm{1_A} = \norm{1_A \cdot 1_A} \le \norm{1_A} \cdot \norm{1_A} \Rightarrow \norm{1_A} =0 $. Also $1_A=0$ und damit $A= \set{0}$.
% subsection 210 (end)

\subsection[Beispiele für normierte Algebren]{Beispiele} % (fold)
\label{sub:211}
\begin{enumerate}[(i)]
	\item $C(\Omega,\mathds{K})$, wobei $\Omega$ ein kompakter Hausdorffraum und die Multiplikation punktweise ist. Betrachte $\norm[\infty]{.}$. Für $f,g : \Omega \to K$ gilt dann
	$\norm[\infty]{f \cdot g} \le \norm[\infty]{f} \cdot \norm[\infty]{g}$
	\item $\mathcal{L}(X,X)$, für $X$ ein normierter Raum, mit $\norm{.}$. Für $S,T \in \mathcal{L}(X,X)$ gilt
	\[
		\norm{S T} = \sup_{\norm[X]{x}=1} \norm[X]{S(T(x))} \le \sup_{\norm[X]{x}=1} \enbrace[\big]{\norm{S} \cdot \norm[X]{Tx}} = \norm{S} \cdot \norm{T}.     
	\]
\end{enumerate}
% subsection 211 (end)

\subsection[Proposition: Norm auf dem Produkt normierter Räume]{Proposition} % (fold)
\label{sub:212}
Seien $X,Y$ normierte Räume. Dann ist $X \times Y$ ein normierter Raum mit 
\begin{align*}
	 \norm[1]{(x,y)} &:= \norm[X]{x} + \norm[Y]{y}   \quad \text{ für } (x,y) \in X \times Y \\
	\text{oder } \quad \norm[\infty]{(x,y)} &:= \max \set{\norm[X]{x}, \norm[Y]{y}}  
\end{align*}
Beide Normen sind äquivalent.
% subsection 212 (end)

\subsection[Proposition und Definition: Norm auf dem Quotientenraum]{Proposition und Definition} % (fold)
\label{sub:213}
Sei $X$ ein normierter Raum und $Y \subseteq X$ ein abgeschlossener Unterraum. Dann ist $\nicefrac{X}{Y} = \set[x +Y]{x \in X}$ ein normierter Raum mit
\[
	\norm[\nicefrac{X}{Y}]{x +Y} := \inf \set[{\norm[X]{x+y} }]{y \in Y}  
\]
Die Quotientenabbildung $q : X \twoheadrightarrow \nicefrac{X}{Y}$, $x \mapsto x+ Y$ ist stetig, linear und $\norm{q} \le 1$ und offen, d.h. bildet offene Mengen in $X$ auch offene Mengen in $\nicefrac{X}{Y}$ ab.
\minisec{Bemerkung}
Daraus folgt, dass Normtopologie und Quotiententopologie auf $\nicefrac{X}{Y}$ gleich sind.
\minisec{Beweis}
\begin{enumerate}[a)]
	\item $\nicefrac{X}{Y}$ ist ein Vektorraum mit $(x+Y) + (x'+Y) = (x+x') + Y$.
	\item Sei $\overline{x} = x+ Y \in \nicefrac{X}{Y}$. Es gilt
	\begin{itemize}
		\item $\norm{\overline{x}} \ge 0 $ Klar für alle $\overline{x} \in \nicefrac{X}{Y}$
		\item Angenommen $\norm[\nicefrac{X}{Y}]{\overline{x}} =0$, d.h. $\inf_{y \in Y} \norm[X]{x+y} =0$. Also existiert eine Folge $(y_n)_n \subset Y$ mit
		$\norm[X]{x+y_n} \xrightarrow{n \to \infty} 0 \Rightarrow y_n \xrightarrow{n \to \infty} -x$. Also ist $-y$ und damit auch $y$ in $Y$, da $Y$ abgeschlossen ist.
		Damit ist $\overline{x}=0 $ 
	\end{itemize}
	\item Sei $0 \not= \lambda \in \mathds{K}$, $x \in X$. Dann gilt
	\begin{align*}
		\norm[\nicefrac{X}{Y}]{\lambda \cdot \overline{x}} = \inf_{y \in Y} \norm[X]{\lambda x + y} = \inf_{y \in Y} \norm[X]{\lambda x + \lambda y} 
		= \inf_{y \in Y} \abs*{\lambda } \cdot \norm[X]{x+y}  =   \abs*{\lambda } \cdot \norm[\nicefrac{X}{Y}]{\overline{x}}  
	\end{align*}
	\item Seien $\overline{x}, \overline{y} \in \nicefrac{X}{Y}$. Dann gilt
	\begin{align*}
		\norm[\nicefrac{X}{Y}]{\overline{x} + \overline{y}} = \inf_{z \in Y} \norm[X]{x+y+z} = \inf_{z,z' \in Y} \norm[X]{x+y+z+z'} &\le \inf_{z,z' \in Y}
		\enbrace[\Big]{\norm[X]{x+z} + \norm[X]{y+z'}} \\ &= \norm[\nicefrac{X}{Y}]{\overline{x}} + \norm[\nicefrac{X}{Y}]{\overline{y}}
	\end{align*}
	Damit ist $\norm[\nicefrac{X}{Y}]{.}$ eine Norm auf $\nicefrac{X}{Y}$.
	\item Linearität von $q$ ist klar. Es gilt
	\[
		\norm[\nicefrac{X}{Y}]{\overline{x}} = \inf_{y \in Y} \norm[X]{x+y} \le \norm[X]{x}   
	\]
	Also
	\[
		\norm{q} = \sup_{\norm[X]{x}=1} \norm[\nicefrac{X}{Y}]{q(x)} \le \sup_{\norm[X]{x}=1 } \norm[X]{x} = 1   
	\]
	\item Zu zeigen: $q$ ist offen. Sei $x \in X$, $\varepsilon>0$. Betrachte die offene Kugel $B(x,\varepsilon) \subset X$. Wir zeigen 
	$B(\overline{x}, \varepsilon) \subset q \enbrace*{B(x,\varepsilon)}$. Sei also $\overline{z} = z + Y$ in $B(\overline{x},\varepsilon)$. Dann gilt
	$\norm[\nicefrac{X}{Y}]{\overline{x}- \overline{z}} < \varepsilon$, worauf folgt
	\[
		\inf_{y \in Y} \norm[X]{x-z+y} < \varepsilon 
	\]
	$\Rightarrow$ es existiert $y \in Y : \norm[X]{x-z+y}<\varepsilon $. Es gilt $\overline{z}=\overline{z-y}\in q\enbrace*{B(x,\varepsilon)}$, da $z-y\in B(x,\varepsilon)$.
	
	Sei $V \subset X$ offen, zeige $q(V) \subset \nicefrac{X}{Y}$ offen. Für $x \in V$ finde $\varepsilon>0$, sodass $B(x,\varepsilon) \subset V$. Dann folgt
	$B(\overline{x},\varepsilon) \subseteq q \enbrace[\big]{ B(x,\varepsilon)} \subseteq q(V)$. \bewende
\end{enumerate}
% subsection 213 (end)

\subsection[Definition: Banachraum und Banachalgebra]{Definition} % (fold)
\label{sub:214}
Ein \Index{Banachraum} ist ein vollständiger, normierter Raum. Eine \Index{Banachalgebra} ist eine vollständige normierte Algebra.  
% subsection 214 (end)

\subsection[Beispiel für eine Banachalgebra]{Beispiel} % (fold)
\label{sub:215}
$C(\Omega, \mathds{K} )$, wobei $\Omega$ ein kompakter Hausdorffraum ist, ist eine Banachalgebra.
% subsection 215 (end)

\subsection[Proposition: Produkte und Quotienten von Banachräumen]{Proposition} % (fold)
\label{sub:216}
Produkte und Quotienten (nach abgeschlossenen Unterräumen) von Banachräumen sind wieder Banachräume.
\minisec{Beweis}
\begin{description}
	\item[Produkt $X \times Y$:] Es gilt $\norm[\infty]{(x,y)} = \max \set{\norm{x}, \norm{y}}$. Sei $\enbrace*{(x_n, y_n)}_n \subset X \times Y$ eine Cauchyfolge.
	Dann sind $(x_n)_n$ und $(y_n)_n$ Cauchyfolgen und es gilt $x_n \to x$ und $y_n \to y$ und somit auch $(x_n, y_n) \to (x,y)$.
	\item[Quotient $\nicefrac{X}{Y}$:] Sei $X$ ein Banachraum und $Y \subseteq X$ ein abgeschlossener Unterraum. 
	
	Falls $\overline{x}, \overline{y} \in \nicefrac{X}{Y}$, so existiert $y' \in X$, sodass $\overline{y} = \overline{y'}$ und 
	\[
		\norm[X]{x-y'} \le 2 \cdot \norm[\nicefrac{X}{Y}]{\overline{x} - \overline{y}}  
	\] 
	Sei nun $\enbrace*{\overline{x}_n }_n$ eine Cauchyfolge in $\nicefrac{X}{Y}$. Es gibt eine Teilfolge $\enbrace*{\overline{x}_{n_k}}_k$, sodass
	\[
		\sum_{k=1}^{\infty} \norm[\nicefrac{X}{Y}]{\overline{x}_{n_{k-1}} - \overline{x}_{n_k}} < \infty 
	\]
	Falls $\enbrace*{\overline{x}_{n_k} }_k $ konvergiert, dann auch $\enbrace*{\overline{x}_n}_n$. Wir dürfen also \oE annehmen, dass $\sum_{n=1}^{\infty} \norm{\overline{x}_{n-1} - \overline{x}_n }< \infty$. Wähle $x_1' \in \overline{x}_1$. Wähle induktiv $x_n' \in X$  mit
	\[
		\norm[X]{x_{n-1}' - x_n'} \le 2 \cdot \norm[\nicefrac{X}{Y}]{\overline{x_{n-1} - \overline{x}_n } }  
	\]
	und $x_n' \in \overline{x}_n$. Dann gilt $\sum_{n=1}^{\infty} \norm[X]{x_{n-1}' - x_n'} < \infty $. Daher ist $(x_n')_n$ eine Cauchyfolge in $X$. Also existiert 
	$x \in X$ mit $x_n \to x$. Dann $\overline{x}_n  \to \overline{x}$. \bewende
\end{description}
% subsection 216 (end)

\subsection[Definition: Sublineare Abbildung]{Definition} % (fold)
\label{sub:217}
Sei $X$ ein $\mathds{R}$-Vektorraum. Eine Abbildung $\varphi : X \to \mathds{R}$ heißt \Index{sublinear}, falls gilt:
\begin{enumerate}[(i)]
	\item $\varphi(\lambda \cdot x) = \lambda \cdot \varphi(x)$, für $x \in X$, $\lambda \in [0,\infty)$
	\item $\varphi(x+y) \le \varphi(x) + \varphi(y)$, für $x,y \in X$.
\end{enumerate}
Setze $\mathcal{S}(X) := \set[\varphi : X \to \mathds{R}]{\varphi \text{ sublinear}}$. $\mathcal{S}(X)$ ist partiell geordnet mit 
\[
	\varphi \le \psi :\iff \varphi(x) \le \psi(x) ,\quad  x \in X
\]
% subsection 217 (end)

\subsection[Beispiele für sublineare Abbildungen]{Beispiele} % (fold)
\label{sub:218}
\begin{enumerate}[(i)]
	\item Halbnormen sind sublinear.
	\item \[
		\begin{tikzpicture}[scale=0.5]
			\draw (-3,0) -- (3,0) node[below right]{$X$};
			\draw (0,-3) -- (0,3) node[right]{$\mathds{R}$};
			\draw[very thick, dashed,OrangeRed2] (0,0) -- (220:4);
			\draw[very thick,SeaGreen4] (40:4) -- (0,0) -- (205:4);
		\end{tikzpicture}
	\]
	$0 = \varphi(0) = \varphi(x-x) =\le \varphi(x) + \varphi(-x)$.
\end{enumerate}
% subsection 218 (end)

\subsection[Proposition: Die sublinearen Abbildungen $\mathcal{S}(X)$ sind nach unten induktiv geordnet]{Proposition} % (fold)
\label{sub:219}
Sei $X$ ein $\mathds{R}$-Vektorraum. $\mathcal{S}(X)$ ist nach unten induktiv geordnet, d.h. jede nichtleere total geordnete Teilmenge besitzt eine untere Schranke.
\minisec{Beweis}
Sei $\emptyset \not= (\varphi_i)_{I} \subseteq \mathcal{S}(X)$ total geordnet. Setze $\varphi(x) := \inf_{i \in I} \varphi_i(x)$, dann gilt 
\[
	- \varphi_i(-x) \le \varphi_i(x) \enspace\Longrightarrow \enspace- \varphi(x) \le \varphi(x) \le \varphi_i(x) < \infty
\]
ebenso $-\varphi(x) \le \varphi(-x) \le \varphi_i(-x)$ für $x \in X$, $i \in I$. Insbesondere gilt $- \infty <  \varphi(x) < \infty$, $x \in X$. 
Die Sublinearität von $\varphi$ und $\varphi \le \varphi_i$, $i \in I$ sind klar (Warum?). \bewende
% subsection 219 (end)

\subsection{Proposition} % (fold)
\label{sub:220}
Sei $X$ ein $\mathds{R}$-Vektorraum und $\varphi \in \mathcal{S}(X)$. Dann ist $\varphi$ in $\mathcal{S}(X)$ minimal genau dann, wenn $\varphi$ linear ist.
\minisec{Beweis}
\begin{description}
	\item["$\Leftarrow$":] Sei $\psi \le \varphi$ mit $\psi$ sublinear und $\varphi$ linear. Dann folgt $\psi(x) \le \varphi(x)$ und $\psi(-x) \le \varphi(-x)$ für 
	$x \in X$.
	\[
		 \Longrightarrow \enspace- \psi(x) \le \psi(-x) \le \varphi(-x) = - \varphi(x)
	\]
	Daraus folgt $\varphi(x) \le \psi(x)$ und somit muss $\varphi=\psi$ gelten. Also ist $\varphi$ minimal.
	\item["$\Rightarrow $":] Sei $\varphi \in \mathcal{S}(X)$ minimal. Zu $\overline{x} \in X$ definiere $\varphi_{\overline{x}} : X \to \mathds{R}$ durch
	\[
		\varphi_{\overline{x}}(x) := \inf_{\lambda \ge 0} \enbrace[\Big]{\varphi \enbrace*{x+ \lambda \cdot \overline{x} } - \lambda \cdot \varphi(\overline{x})}
	\]
	Es gilt $\varphi_{\overline{x}}(x) \in (-\infty,\infty)$, denn für $x \in X, \lambda \ge 0$ ist \marginnote{mit (ii) aus \ref{sub:217}}
	\[
		-\varphi(-x) \le \varphi(x+ \lambda \cdot \overline{x}) - \lambda  \cdot \varphi(\overline{x}) \le \varphi(x)
	\]
	$\varphi_{\overline{x}}$ ist sublinear:
	\begin{enumerate}[(i)]
		\item Sei $\mu >0$. Dann gilt für $x \in X$.
		\begin{align*}
			\varphi_{\overline{x}}(\mu \cdot x) = \inf_{\lambda \ge 0} \enbrace[\Big]{ \varphi \enbrace*{\mu \cdot x + \lambda  \cdot \overline{x} } 
			- \lambda \cdot \varphi(\overline{x} )} &= \inf_{\lambda \ge 0 } \mu \cdot  \enbrace*{ \varphi \enbrace*{x + \frac{\lambda }{\mu} \cdot \overline{x}}
			- \frac{\lambda }{\mu} \cdot \varphi(\overline{x} )  } \\
			&= \mu \cdot \inf_{\lambda' \ge 0} \enbrace[\Big]{ \varphi \enbrace*{x+ \lambda' \cdot \overline{x} } - \lambda' \varphi(\overline{x})} \\
			&= \mu \cdot \varphi_{\overline{x}}(x) 
		\end{align*}
		$\varphi_{\overline{x}}(0 \cdot x) = 0$ ist klar.
		\item Zu $x,y \in X$, $\varepsilon>0$ wähle $\lambda_x, \lambda_y \ge 0$ mit 
		\begin{align*}
			\varphi_{\overline{x}}(x) &\ge \varphi \enbrace*{x + \lambda_x \cdot \overline{x}} - \lambda_x \cdot \varphi(\overline{x}) - \varepsilon \\
			\varphi_{\overline{x}}(y) &\ge \varphi \enbrace*{y + \lambda_y \cdot \overline{x}} - \lambda_y \cdot \varphi(\overline{x}) - \varepsilon 
		\end{align*}
		Setze $\lambda  := \lambda_x + \lambda_y$, dann gilt
		\begin{align*}
			\varphi_{\overline{x}}(x) + \varphi_{\overline{x}}(y) &\ge \varphi \enbrace*{ x+ \lambda_x \cdot \overline{x}} + 
			\varphi \enbrace*{y+ \lambda_y \cdot \overline{x} } - \lambda  \cdot \varphi(\overline{x}) - 2 \varepsilon \\
			&\ge  \varphi \enbrace*{x +y + \lambda \cdot \overline{x}} - \lambda  \cdot \varphi(\overline{x}) - 2 \varepsilon \\
			&\ge \varphi_{\overline{x}} (x+y) - 2 \varepsilon
		\end{align*}
		Da $\varepsilon>0$ beliebig war, gilt $\varphi_{\overline{x}}(x) + \varphi_{\overline{x}}(y) \ge \varphi_{\overline{x}}(x+y)$. Also ist $\varphi_{\overline{x}}$
		sublinear.
	\end{enumerate}
	$\varphi_{\overline{x}} \le \varphi$ ist klar. Da $\varphi$ minimal ist, folgt $\varphi_{\overline{x}} = \varphi$. Wir erhalten 
	\[
		\varphi(x) + 1 \cdot \varphi(\overline{x}) \le \varphi(x+ 1 \cdot \overline{x}) \le \varphi(x) + 1 \cdot \varphi(\overline{x})
	\]
	für jedes $x, \overline{x}\in X$. Also ist $\varphi$ additiv und somit linear. \bewende
\end{description}
% subsection 220 (end)

\subsection{Satz von Hahn-Banach} % (fold)
\label{sub:221}
Sei $X$ ein $\mathds{R}$-Vektorraum und $\varphi \in \mathcal{S}(X)$.\marginnote{$X$ ist nicht zwangsweise endlich-dimensional!}
Dann existiert $\psi : X \to \mathds{R}$ linear mit $\psi \le \varphi$.
\minisec{Beweis}
Sei $\mathcal{S}_\varphi : =\set[\varphi' \in \mathcal{S}(X)]{\varphi' \le \varphi} \ni \varphi$. Nach Proposition \ref{sub:219} ist $\mathcal{S}(X)$ und damit auch 
$\mathcal{S}_\varphi$ nach unten induktiv geordnet. Nach dem Lemma von Zorn enthält $\mathcal{S}_\varphi$ ein minimales Element $\psi$. $\psi$ ist auch minimal in 
$\mathcal{S}(X)$ \hfill (Warum?)\\
Also ist $\psi$ linear nach Proposition \ref{sub:220}. \bewende \medskip \\
Wichtige Folgerungen: Fortsetzungs- und Trennungssätze.
% subsection 221 (end)

\subsection[Satz (Hahn-Banach): Existenz einer linearen Fortsetzung]{Satz} % (fold)
\label{sub:222}
Sei $X$ ein $\mathds{R}$-Vektorraum und $\varphi : X \to \mathds{R}$ sublinear. Sei $Y \subset X$ ein linearer Unterraum und $\psi : Y \to \mathds{R}$ linear mit
$\psi \le \varphi\big|_{Y}$. Dann existiert $\overline{\psi} : X \to \mathds{R}$ linear mit $\overline{\psi}\big|_Y = \psi$ und $\overline{\psi} \le \varphi$.
\minisec{Beweis}
Definiere $\tilde{\varphi} : X \to \mathds{R}$ durch $\tilde{\varphi}(x) := \inf_{y \in Y} \enbrace[\big]{ \varphi(x-y) + \psi(y)}$, $x \in X$. Es gilt
\[
	\varphi(x-y) + \psi(y) \ge \varphi(-y) - \varphi(-x) - \psi(-y) \ge - \varphi(-x) > -\infty \qquad \forall x \in X, y \in Y
\]
Also ist $\tilde{\varphi}$ wohldefiniert. $\tilde{\varphi}$ ist sublinear: vergleiche \ref{sub:220}. Nach Satz von Hahn-Banach (\ref{sub:221}) existiert ein 
$\overline{\psi} : X \to \mathds{R}$ linear mit $\overline{\psi} \le \tilde{\varphi}$. Aus \ref{sub:220} folgt, dass $\overline{\psi}$ minimal in $\mathcal{S}(X)$ ist.
Ebenso ist $\overline{\psi}\big|_{Y}$ minimal in $\mathcal{S}(Y)$. Also gilt
\[
	\overline{\psi} \le \tilde{\varphi}\big|_{Y} \le \psi \quad \Rightarrow \quad \psi = \overline{\psi}\big|_{Y} \bewende  
\]
% subsection 222 (end)

\subsection{Satz} % (fold)
\label{sub:223}
Sei nun $X$ ein $\mathds{K}$-Vektorraum und $p : X \to \mathds{R}$ eine Halbnorm. Sei $Y \subset X$ ein Untervektorraum und $\psi : Y \to \mathds{K}$ linear mit
$\abs*{\psi(y)} \le p(y)$ für $y \in Y$. Dann existiert $\overline{\psi} : X \to \mathds{K}$ linear mit $\overline{\psi}\big|_{Y} = \psi$ und 
\[
	\abs*{\overline{\psi}(x) } \le p(x) \quad \text{ für } x \in X 
\]
\minisec{Beweis}
Sei zunächst $\mathds{K}=\mathds{R}$. $p$ ist sublinear und es gilt $\psi \le p\big|_{Y}$. Nach \ref{sub:222} existiert eine lineare Fortsetzung 
$\overline{\psi} : X \to \mathds{R}$ mit $\overline{\psi} \le p$. Es gilt auch
\begin{align*}
	- \overline{\psi}(x) =  \overline{\psi}(-x) \le p(-x)  = p(x) \quad \Longrightarrow \quad \abs*{\overline{\psi}(x) } \le p(x) \quad \text{ für } x \in X
\end{align*}
Sei nun $\mathds{K}=\mathds{C}$. Definiere $\psi_1 := \re(\psi) : Y \to \mathds{R}$. $\psi_1$ ist $\mathds{R}$-linear (warum?). Es gilt 
$\abs*{\psi_1(y)} \le \abs*{\psi(y)} \le p(y)$, $y \in Y$. Es existiert also ein $\overline{\psi}_1 : X \to \mathds{R}$ $\mathds{R}$-linear mit
\[
	\overline{\psi}_1 (y) = \psi_1(y), \enspace y \in Y \quad \text{ und } \quad \abs*{\overline{\psi}_1(x)} \le p(x) , \enspace x \in X 
\]
Definiere jetzt $\overline{\psi} : X \to \mathds{C}$ durch $\overline{\psi}(x) := \overline{\psi}_1(x) - i \cdot \overline{\psi}_1(i \cdot x)$, $x \in X$.
Dann ist $\overline{\psi}$ $\mathds{C}$-linear (warum?). Weiter gilt $\overline{\psi}(y) = \psi(y)$ für $y \in Y$.
($\re(\overline{\psi}\big|_{Y})= \re(\psi)$ und $\overline{\psi}\big|_{Y}$ und $\psi$ sind beide $\mathds{C}$-linear).

Zu $x \in X$ wähle $\lambda \in \mathds{C}$ mit $\abs*{\lambda }=1$ und 
\[
	\abs*{\overline{\psi}(x)} = \lambda  \cdot \overline{\psi}(x) = \overline{\psi}(\lambda \cdot x) = \overline{\psi}_1(\lambda \cdot x) \le p(\lambda \cdot x)
	= \abs*{\lambda } \cdot p(x) = p(x) \bewende     
\]
% subsection 223 (end)

\subsection{Satz} % (fold)
\label{sub:224}
Sei $X$ ein normierter $\mathds{K}$-Vektorraum, $Y \subset X$ ein Unterraum und $\psi : Y \to \mathds{K}$ linear und stetig. Dann existiert eine lineare Fortsetzung 
$\overline{\psi} : X \to \mathds{K}$ mit $\norm[\mathcal{L}(X,\mathds{K})]{\overline{\psi}} = \norm[\mathcal{L}(Y,\mathds{K})]{\psi}$.
\minisec{Beweis}
Definiere eine Halbnorm durch $p(x) := \norm{x} \cdot \norm{\psi}$. Nach Satz \ref{sub:223} besitzt $\psi$ eine lineare Fortsetzung $\overline{\psi}$ mit
$\abs*{\overline{\psi(x)}} \le p(x)= \norm{\psi} \cdot \norm{x}$, $x \in X$. Es folgt $\norm{\overline{\psi}} \le \norm{\psi}$. $\norm{\psi} \le \norm{\overline{\psi}}$
ist trivial, da $\psi\big|_{Y} = \overline{\psi}$. \bewende
% subsection 224 (end)

\subsection[Definition: Konvexe Teilmenge eines $\mathds{K}$-Vektorraums]{Definition} % (fold)
\label{sub:225}
Sei $X$ ein $\mathds{K}$-Vektorraum. Sei $M \subset X$ eine Teilmenge. $M$ heißt \Index{konvex}, falls für $a,b \in M$, $\lambda \in [0,1]$ gilt
\[
	(1-\lambda ) \cdot a + \lambda \cdot b \in M
\]
Ist $X$ ein topologischer Vektorraum und $M \subset X$ konvex, so ist auch $\overline{M}$ konvex. \hfill (warum?)\\
Ist $N \subset X$ eine beliebige Teilmenge, so ist 
\[
	\conv(N) := \smashoperator{\bigcap_{\substack{N \subset M \subset X \\ M \text{ konvex} }}} \, M
\]
$\conv(N)$ ist konvex. (warum?)
% subsection 225 (end)

\subsection{Satz} % (fold)
\label{sub:226}
Sei $X$ ein $\mathds{R}$-Vektorraum, $\emptyset \not= M \subset X$ konvex, $\varphi \in \mathcal{S}(X)$. Dann existiert $\psi : X \to \mathds{R}$ linear mit 
$\psi \le \varphi$ und 
\begin{equation*}
	\inf_{y \in M} \varphi(y) = \inf_{y \in M} \psi(y) \label{eq:226:1}\tag{*}
\end{equation*}
\minisec{Beweis}
Setze $\mu := \inf_{y \in M} \varphi(y)$. Falls $\mu = -\infty$, so folgt die Behauptung aus dem Satz von Hahn-Banach \ref{sub:221}. $\eqref{eq:226:1}$ ist trivialerweise
erfüllt. Sei also $\mu \in \mathds{R}$. Definiere $\tilde{\varphi} : X \to \mathds{R}$ durch
\[
	\tilde{\varphi}(x) := \inf_{y \in M, \lambda  \ge 0} \enbrace[\big]{ \varphi(x+ \lambda \cdot y) - \lambda  \cdot \mu} , \quad x \in X
\]
Es gilt $\varphi(x + \lambda \cdot y)- \lambda \cdot \mu \ge - \varphi(-x)$, $x \in X$, $y\in M$, $\lambda \ge 0$. Also ist $\tilde{\varphi}(x) \ge - \varphi(-x) >-\infty$.
$\tilde{\varphi}$ ist sublinear:
\begin{enumerate}[(i)]
	\item $\tilde{\varphi}(\gamma \cdot x) = \gamma \cdot \tilde{\varphi}(x)$, für $x \in X$, $\gamma \in \ge 0$, wie in \ref{sub:222}.
	\item Seien $x,z \in X$ und $\varepsilon>0$. Wähle $y_x, y_z \in M$, $\lambda_x, \lambda_y \ge 0$ mit
	\[
		\tilde{\varphi}(x) \ge \varphi (x +\lambda_x \cdot y_x) - \lambda_x \cdot \mu -\varepsilon \quad \text{ und }\quad 
		\tilde{\varphi}(x) \ge \varphi (z +\lambda_z \cdot y_z) - \lambda_z \cdot \mu -\varepsilon
	\] 
	Es folgt 
	\begin{align*}
		\tilde{\varphi}(x) + \tilde{\varphi}(z) &\ge \varphi \enbrace[\big]{x+z+ \lambda_x \cdot y_x + \lambda_z \cdot y_z} - (\lambda_x + \lambda_z) \cdot \mu 
		- 2 \varepsilon \\
		&= \varphi \enbrace[\Bigg]{x+z + (\lambda_x+ \lambda_z)\cdot \underbrace{\enbrace*{\frac{\lambda_x}{\lambda_x + \lambda_z} \cdot y_x + \frac{\lambda_z}{\lambda_x + \lambda_z} \cdot y_z }}_{\in M} } - (\lambda_x + \lambda_z) \cdot \mu -2 \varepsilon \\
		&\ge \tilde{\varphi}(x+z) - 2 \varepsilon
	\end{align*}
\end{enumerate}
Nach Hahn-Banach (\ref{sub:221}) existiert $\psi : X \to \mathds{R}$ linear mit $\psi \le \tilde{\varphi} \le \varphi$. Für $y \in M$ gilt 
\[
	-\psi(y) = \psi(-y) \le \tilde{\varphi}(-y) \le \varphi(-y + 1 \cdot y)- 1 \cdot \mu= -\mu
\]
also $\mu \le \psi(y) \le \varphi(y)$, woraus \eqref{eq:226:1} folgt. \bewende
% subsection 226 (end)

\subsection{Satz} % (fold)
\label{sub:227}
Sei $X$ ein normierter $\mathds{R}$-Vektorraum und $A,B \subset X$ nichtleere konvexe Teilmengen mit \index{Distanz $\dist$}
\[
	\dist(A,B) := \inf \set[\norm{a-b}]{a \in A, b \in B} > 0
\]
Dann existiert $\psi : X \to \mathds{R}$ stetig und linear mit $\psi(A) \cap \psi(B) \not= \emptyset$.
\minisec{Beweis}
$A - B := \set[a-b]{a \in A, b \in B} \subset X$ ist konvex. (Warum?) Nach Satz \ref{sub:226} existiert $\psi : X \to \mathds{R}$ linear mit $\psi \le \norm{.}$ und 
\[
	0 < \dist(A,B)=  \inf_{y \in A -B} \norm{y} = \inf_{y \in A-B} \psi(y) = \inf_{a \in A} \psi(a) - \sup_{b \in B} \psi(b). \bewende
\]
% subsection 227 (end)
% section 2 (end)
\newpage

\section{Operatoren zwischen Banachräumen. Die Sätze von der offenen Abbildung und vom abgeschlossenen Graphen} % (fold)
\label{sec:3}

\subsection[Proposition: $Y$ vollständig $\Rightarrow $ $\mathcal{L}(X,Y)$ ist vollständig]{Proposition} % (fold)
\label{sub:31}
Seien $X,Y$ normierte Vektorräume, $Y$ vollständig. Dann ist $\mathcal{L}(X,Y)$ vollständig bezüglich $\norm[\mathcal{L}(X,Y)]{.}$.
\minisec{Beweis}
Sei $(T_n)_{n \in \mathds{N}} \subseteq \mathcal{L}(X,Y)$ eine Cauchy-Folge bezüglich $\norm[\mathcal{L}(X,Y)]{.}$. Für $x \in X$ ist dann auch $(T_n x)_{n \in \mathds{N}} \subset Y$ Cauchy
bezüglich $\norm[Y]{.}$. (Warum?)\footnote{mit $\norm[Y]{T_n x- T_m x} \le \norm{T_n -T_m} \cdot \norm{x}$}
Da $Y$ ein Banachraum ist, folgt $T_n x \xrightarrow{n \to \infty} y$ für ein $y \in Y$. Wir definieren $T : X \to Y$ durch $T(x) := \lim_{ n \to \infty} T_n(x)$ für 
$x \in X$. $T$ ist linear: Klar (Warum?).

$T$ ist stetig: $(T_n)_{n \in \mathds{N}}$ ist Cauchy, also auch $\enbrace*{\norm{T_n}}_{n \in \mathds{N}} \subset \mathds{R}$ eine Cauchyfolge (Warum?)\footnote{mit 
$\norm{T_n} = \norm{T_n - T_{n_0}+ T_{n_0}} \le \norm{T_n- T_{n_0}} + \norm{T_{n_0}}$}. Dann existiert $C \ge 0$ mit $\norm[\mathcal{L}(X,Y)]{T_n} \le C $, $n \in \mathds{N}$, also
gilt $\norm[Y]{T x} \le C \cdot \norm{x}$, $x \in X$, woraus $\norm[\mathcal{L}(X,Y)]{T} \le C$ und die Stetigkeit von $T$ folgt. Also ist $T \in \mathcal{L}(X,Y)$.

\noindent Falls $\varepsilon >0$, so existiert $n_0 \in \mathds{N}$ mit $\norm[\mathcal{L}(X,Y)]{T_n - T_m}  < \varepsilon$ für $n,m \ge n_0$. Insbesondere: 
\begin{align*}
	\norm[Y]{(T_n- T_m)x} \le \norm[\mathcal{L}(X,Y)]{T_n -T_m} \cdot \norm[X]{x} < \varepsilon \cdot \norm[X]{x}    
\end{align*}
für $n,m \ge n_0$, $x \in X$. Weiter gilt 
\begin{align*}
	\norm[Y]{(T_n- T)x} \le \norm[Y]{(T_n - T_m)x} + \norm[Y]{(T_m-T)x} &\le \norm[\mathcal{L}(X,Y)]{T_n - T_m} \cdot \norm[X]{x} + \norm[Y]{(T_m-T)x} \\
	&\le \varepsilon \cdot \norm[X]{x} + \varepsilon \cdot \norm[X]{x}  
\end{align*}
für $n,m \ge n_0$ und $m$ groß genug, $x \in X$. Damit folgt nun
$\norm[\mathcal{L}(X,Y)]{T_n -T} \le \varepsilon$ für$n \ge n_0$, also $T_n \to T$ bezüglich $\norm[\mathcal{L}(X,Y)]{.}$. \bewende
% subsection 31 (end)

\subsection[Corollar: Der Dualraum ist vollständig, $\mathcal{L}(X,X)$ ist Banachalgebra, falls $X$ Banachraum ist]{Corollar} % (fold)
\label{sub:32}
\begin{enumerate}[(i)]
	\item Der Dualraum $X^*$ eines normierten $\mathds{K}$-Vektorraumes $X$ ist vollständig. 
	\item Falls $\set{0} \not= X$ ein Banachraum ist, so ist $\mathcal{L}(X,X)$ eine Banachalgebra.
\end{enumerate}
\minisec{Beweis}
\begin{enumerate}[(i)]
	\item $X^*=\mathcal{L}(X,\mathds{K})$ und $\mathds{K}$ ist vollständig. Wende \ref{sub:31} an.
	\item $\mathcal{L}(X,X)$ ist ein Banachraum nach \ref{sub:31}. $\mathcal{L}(X,X)$ ist eine normierte Algebra: 
	\begin{itemize}
		\item $\norm[\mathcal{L}(X,X)]{\id_X}=1$ \hfill (Warum?)
		\item Für $S,T \in \mathcal{L}(X,X)$, $x \in X$ gilt
		\begin{align*}
			\norm[X]{(S \circ T)x} \le \norm{S} \cdot \norm[X]{T x} \le \norm{S} \cdot \norm{T} \cdot \norm[X]{x}   
		\end{align*}
		Also ist $\norm{S \circ T} \le \norm{S} \cdot \norm{T}$. \bewende
	\end{itemize}
\end{enumerate}
% subsection 32 (end)

\subsection[Definition und Proposition: lineare Isometrie $\iota_X : X \to X^{**}$]{Definition und Proposition} % (fold)
\label{sub:33}
Sei $X$ ein normierter Raum. Definiere $\iota_X : X \to X^{**} \enspace{\color{light_gray}\enbrace[\big]{=(X^*)^*}}$ durch
\[
	x \longmapsto \enbrace[\big]{\varphi \mapsto \varphi(x) } 
\]
für $x \in X$. Die Abbildung $\iota_X$ ist eine lineare Isometrie.
\minisec{Beweis}
\begin{description}
	\item[Wohldefiniertheit:] Die Linearität von $\iota_X(x)$ ist klar. Mit
	\begin{align*}
		\abs*{\iota_X(x)(\varphi)} = \abs*{\varphi(x)} \le \norm[X^*]{\varphi} \cdot \norm[X]{x} = \norm[X]{x} \cdot \norm[X^*]{\varphi}      
	\end{align*}
	folgt $\varphi \in X^*$ und somit $\norm[X^{**}]{\iota_X(x)} \le \norm[X]{x}$, also ist $\iota_X(x)$ stetig.
	\item[$\iota_X$ ist linear:] Für $\alpha, \beta \in \mathds{K}$, $x,y \in X$, $\varphi \in \mathcal{L}(X,\mathds{K})= X^*$ gilt
	\begin{align*}
		\iota_X(\alpha \cdot x + \beta \cdot y)(\varphi)= \varphi(\alpha \cdot x + \beta \cdot y) = \alpha \circ \varphi(x) + \beta \cdot \varphi(y) 
		&= \alpha \cdot \iota_X(x)(\varphi) + \beta \cdot \iota_X(y)(\varphi) \\
		&= \enbrace[\big]{\alpha \cdot \iota_X(x) + \beta \cdot \iota_X(y)}(\varphi) 
	\end{align*}
	\item[Isometrie:] Es bleibt zu zeigen: $\norm[X^{**}]{\iota_X(x)} \ge \norm[X]{x}$ für $0 \not= x \in X$. Sei $Y := \mathds{K}\cdot x$ und $\psi : Y \to \mathds{K}$
	gegeben durch $\psi(\alpha \cdot x) := \alpha \cdot \norm[X]{x}$. Dann ist $Y \subset X$ ein linearer Unterraum und $\psi(Y) \to \mathds{K}$ linear mit
	$\norm{\psi}=1$. Nach Satz \ref{sub:224} existiert $\overline{\psi}. X \to \mathds{K}$ linear mit $\overline{\psi}(x) = \psi(x) = \norm[X]{x}$ und 
	$\norm{\overline{\psi}}= \norm{\psi} = 1$. Es gilt 
	\begin{align*}
		\norm[X^{**}]{\iota_X(x)} \ge \abs*{\iota_X(x)(\overline{\psi})} = \abs*{\overline{\psi}(x)} = \norm[X]{x}. \bewende
	\end{align*}
\end{description}
% subsection 33 (end)

\subsection{Definition und Proposition} % (fold)
\label{sub:34}
Seien $X,Y$ normierte Räume. Definiere eine Abbildung $.^{\tr} : \mathcal{L}(X,Y) \to \mathcal{L}(Y^*,X^*)$ durch $T \mapsto T^{\tr} = (\varphi \mapsto \varphi \circ T)$.
$.^\tr$ ist eine lineare Isometrie und das Diagramm
\[
	\begin{tikzcd}
		X \rar["\iota_X"] \dar["T"]& X^{**} \dar["T^{\tr \, \tr}= (T^\tr)^\tr"]\\
		Y \rar["\iota_Y"]& Y^{**}
	\end{tikzcd}
\]
kommutiert für jedes $T \in \mathcal{L}(X,Y)$. Für $\iota_{X^*} : X^* \to X^{***}$ und $\iota_{X}^\tr : X^{***} \to X^*$ gilt $\iota_X^\tr \circ \iota_{X^*} = \id_{X^*}$.
\minisec{Beweis}
\begin{enumerate}[(i)]
	\item Für $T,S \in \mathcal{L}(X,Y)$, $\alpha, \beta \in \mathds{K}$, $\varphi \in Y^*$ gilt
	\begin{align*}
		\enbrace*{\alpha \cdot T + \beta \cdot S}^\tr (\varphi) = \varphi \circ \enbrace*{\alpha \cdot T + \beta \cdot S} \!\!\stackrel{\text{linear}}{=} \!\!
		\alpha \cdot (\varphi \circ T) + \beta \cdot (\varphi \circ S)  &= \alpha \cdot T^\tr(\varphi) + \beta \cdot S^\tr(\varphi)\\
		&= \enbrace*{\alpha \cdot T^\tr + \beta \cdot S^\tr}(\varphi) 
	\end{align*}
	\item Für $x \in X$, $\psi \in Y^*$ gilt $\enbrace[\big]{\iota_Y T(x)}(\psi) = \psi \enbrace[\big]{T(x)}$; weiter gilt
	\begin{align*}
		\enbrace[\big]{T^{\tr \, \tr} \iota_X(x)} (\psi) = \enbrace[\big]{(T^\tr)^\tr \iota_X(x)}(\psi)  = \enbrace[\big]{\iota_X(x) \circ T\tr}(\psi)
		&= \iota_X(x) \enbrace*{T^\tr(\psi)}  \\
		&= T^\tr(\psi)(x) = \psi \circ T(x) \\
		&= \iota_Y \enbrace[\big]{T(x)} (\psi) 
	\end{align*}
	$\Rightarrow T^{\tr \, \tr} \iota_X(x) = \iota_Y \enbrace[\big]{T(x)} = \iota_Y \circ T (x)$. Also $T^{\tr \, \tr} \circ \iota_X = \iota_Y  \circ T$. Also kommutiert
	das Diagramm.
	\item Es ist $\abs{(T^\tr \psi)(x)} = \abs*{\psi T(x)} \le \norm[Y^*]{\psi} \cdot \norm[Y]{T x} \le \norm[Y^*]{\psi} \cdot \norm[\mathcal{L}(X,Y)]{T} \cdot \norm[X]{x}$. Also
	ist $\norm[X^*]{T^\tr \psi} \le  \norm[Y^*]{\psi} \cdot  \norm[\mathcal{L}(X,Y)]{T}$ und damit folgt 
	\[
		\norm[\mathcal{L}(Y^*,X^*)]{T^\tr} \le \norm[\mathcal{L}(X,Y)]{T}  
	\]
	Ebenso ist $\norm[\mathcal{L}(X^{**},Y^{**})]{T^{\tr \, \tr}} \le \norm[\mathcal{L}(Y^*,X^*)]{T^\tr} \le \norm[\mathcal{L}(X,Y)]{T}$. Andererseits gilt 
	$\norm[l(X,Y)]{T} \le \norm[\mathcal{L}(X^{**},Y^{**})]{T^{\tr\, \tr}}$ nach \ref{sub:33} und da das Diagramm kommutiert (Warum?). 
	\[
		\Rightarrow \norm[\mathcal{L}(X,Y)]{T} \le \norm[\mathcal{L}(Y^*,X^*)]{T^\tr}  \le  \norm[\mathcal{L}(X,Y)]{T} 
	\]
	also ist $.^\tr$ eine Isometrie.
	\item Für $\varphi \in X^*$, $x \in X$ gilt
	\begin{align*}
		\enbrace[\big]{\iota_X^\tr \circ \iota_{X^*} (\varphi)} (x) = \iota_{X^*}(\varphi) \circ \iota_X (x) = \iota_X(x)(\varphi) = \varphi(x)
	\end{align*}
	$\Rightarrow \iota_X^\tr \circ \iota_{X^*}(\varphi)= \varphi$, also $\iota_X^\tr \circ  \iota_{X^*} = \id_{X^*}$. \bewende
\end{enumerate}
% subsection 34 (end)

\subsection{Satz: Prinzip der gleichmäßigen Beschränktheit für Banachräume} % (fold)
\label{sub:35}
Sei $X$ ein Banachraum und $Y$ ein normierter Raum. Sei $M \subset \mathcal{L}(X,Y)$, sodass die Menge 
\[
	\set[{x \mapsto \norm[Y]{T x}}]{T \in M \rule{0cm}{1em}} \subset C(X,\mathds{R} )
\]
punktweise gleichmäßig beschränkt ist, d.h. für $x \in X$ existiert $C_x \ge 0$ mit $\norm[Y]{T x} > C_x$, $T \in M$. Dann existiert $C \ge 0$ mit $\norm[\mathcal{L}(X,Y)]{T} \le C$,
$T \in M$.
\minisec{Beweis}
Nach Corollar \ref{sub:112} (Prinzip der gleichmäßigen Beschränktheit für metrische Räume) existieren eine offene Kugel $\emptyset \not= B(x_0,\varepsilon) \subset X$ und 
$K \ge 0$ mit $\norm[Y]{T x} \le K$ für $x \in B(x_0, \varepsilon)$, $T \in M$. Für $x \in X$ mit $\norm{x} \le 1$ gilt 
\begin{align*}
	\norm[X]{T x} = \frac{e}{\varepsilon} \cdot  \norm[Y]{T \enbrace*{\frac{\varepsilon}{2} \cdot x}} &= \frac{2}{\varepsilon} \cdot 
	\norm[Y]{T \enbrace*{\frac{\varepsilon}{2} \cdot x + x_0 - x_0 } }  \\
	&\le \frac{2}{\varepsilon} \cdot   \enbrace*{ \doppelstrich[\Bigg]{T \enbrace[\Bigg]{\underbrace{\frac{\varepsilon}{2} \cdot x + x_0}_{\in B(x_0,\varepsilon)} } }_Y +
	 \norm[Y]{T  \enbrace[\Big]{\underset{\in B(x_0,\varepsilon)}{x_0}} }} \\
	&\le \frac{2}{\varepsilon} (K+K) = \frac{4K}{\varepsilon}   
\end{align*}
$\Rightarrow \norm[\mathcal{L}(X,Y)]{T} \le \frac{4 K}{\varepsilon} =: C$. \bewende
% subsection 35 (end)

\subsection{Corollar} % (fold)
\label{sub:36}
Sei $Z$ ein metrischer Raum und $N \subset Z$ eine Teilmenge, so dass gilt: Für alle $\varphi \in Z^*$ existiert $C_\varphi \ge 0$ mit $\abs*{\varphi(z)} \le C_\varphi$, 
$z \in N$.  Dann ist $N$ beschränkt.
\minisec{Beweis}
$\iota_Z : Z \to Z^{**}$ ist eine Isometrie nach \ref{sub:33}. Wende Satz \ref{sub:35} an mit $X=Z^{*}$, $X=\mathds{K}$ und 
$M = \iota_Z(N) \subset Z^{**}= \mathcal{L}(Z^*,\mathds{K})$.
Für $\varphi \in X = Z^{*}$ existiert $C_\varphi \ge 0$ mit $\abs*{T \varphi} = \abs*{\iota_Z(z) (\varphi)} = \abs*{\varphi(z)} \le C_\varphi$, wo $T= \iota_Z(z) \in M$
für ein $z \in N$. Nach \ref{sub:35} folgt: Es existiert ein $C \ge 0$ mit $\norm[Z]{z} = \norm[Z^{**}]{\iota_Z(z)}  \le C$, $z \in N$. \bewende
% subsection 36 (end)

\subsection{Corollar} % (fold)
\label{sub:37}
Sei $X$ ein Banachraum und $Z$ ein normierter Raum. $M \subset \mathcal{L}(X,Z)$ eine Teilmenge, so dass gilt: Für alle $\varphi \in Z^*$ und $x \in X$ existiert $C_{x,\varphi} \ge 0$
mit
\[
	\abs*{\varphi \enbrace[\big]{T(x)}} \le C_{x,\varphi}, \quad T \in M
\]
Dann ist $M$ beschränkt.
\minisec{Beweis}
Für jedes $x \in X$ folgt aus Corollar \ref{sub:36} mit $N := \set[T(x)]{T \in M}$, dass $N$ beschränkt ist. Es folgt, dass $ \set[x \mapsto \norm{T(x)} ]{T \in M}$ ist 
punktweise gleichmäßig beschränkt. Mit Satz \ref{sub:35} folgt: $M$ ist beschränkt. \bewende
% subsection 37 (end)

\subsection{Satz: Prinzip der offenen Abbildung} % (fold)
\label{sub:38}
Seien $X,Y$ Banachräume, $T \in \mathcal{L}(X,Y)$ surjektiv. Dann ist $T$ offen, d.h. für eine offene Menge $U \subset X$ ist $T(U) \subset Y$ offen.
\minisec{Beweis}
\begin{enumerate}[(i)]
	\item Sei $0 \in W \subset X$ offen, dann existiert $\emptyset \not= V \subset Y$ offen mit $V \subset \overline{T(W)}$: $0 \in W \subset X$ $\Rightarrow $ 
	$B(0,\varepsilon) \subset W$ für ein $\varepsilon>0$. Dann gilt
	\[
		X=\bigcup_{n \in \mathds{N}} n \cdot B(0,\varepsilon) = \bigcup_{n \in \mathds{N}} n \cdot W
	\]
	Aus der Surjektivität von $T$ folgt $Y = \bigcup_{n \in \mathds{N} }n \cdot T(W) = \bigcup_{n \in \mathds{N}} n \cdot \overline{T(W)}$. Nach dem Satz von Baire 
	(\ref{sub:110}) folgt: Ein $n \cdot \overline{T(W)}$ enthält eine nichtleere offene Teilmenge, also auch $\overline{T(W)}$.
	\item Sei $= \in M \subset X$ offen, dann ist $0$ ein innerer Punkt von $\overline{T(M)}$: 
	
	Die Abbildung $X \times X \to X$, $(x_1,x_2) \mapsto x_1 -x_2$ ist stetig.\footnote{Banachräume sind topologische Vektorräume} $\Rightarrow$ es gibt $0 \in W \subset X$ 
	offen mit
	\[
		W - W = \set[x_1 -x_2]{x_1,x_2 \in W} \subset M 
	\]
	Es folgt 
	\[
		\overline{T(M)} \supset \overline{T(W-W)} = \overline{T(W)- T(W)} \stackrel{\text{(Warum?)}}{\supset} \overline{T(W)}- \overline{T(W)} \stackrel{\text{(i)}}{\supset}
		V -V
	\]
	für $\emptyset \not= V$ offen in $Y$ wie in (i). Aber $0 \in V-V = \bigcup_{y \in V}\set{y}- V$ ist offen in $Y$. Also ist $0$ ein innerer Punkt von $\overline{T(M)}$.
	\item Ist $0 \in N \subset X$ offen, so existiert $0 \in Z \subset Y$ offen mit $Z \subset T(M)$:
	
	Sei $\varepsilon_0 >0$. Wähle $\varepsilon_i >0$, $1 \le i \in \mathds{N}$ mit $\sum_{i=1}^{\infty} \varepsilon_i < \varepsilon_0$. Nach (ii) gilt: Zu $\varepsilon_i$, 
	$i \in N$ existiert ein $\delta_i >0$ mit $B_Y(0,\delta_i) \subset \overline{T \enbrace*{B_X(0,\varepsilon_i)}}$. Es gilt $\delta_i \xrightarrow{i \to \infty} 0$. Sei 
	nun $y \in B_Y(0,\delta_0) =: Z$. Dann gilt $y \in \overline{T \enbrace*{B_X(0,\varepsilon_0)}}$. $\Rightarrow $ es existiert $x_0 \in B_X(0,\varepsilon)$ mit 
	$\norm[Y]{y- T(x_0)}< \delta_1$. Folglich existiert $x_1 \in B_X(0,\varepsilon_1)$ mit $\norm[Y]{y- T(x_0) - T(x_1)} < \delta_2$. Per Induktion folgt
	\[
		\exists x_i \in B_X(0,\varepsilon_i) \text{ mit } \norm[Y]{y- \sum\nolimits_{j=0}^{i}T(x_i)} < \delta_{i+1}, \qquad i \in \mathds{N} 
	\]
	$\sum_{i \in \mathds{N}}x_i$ konvergiert in $X$ und $\norm[X]{\sum_{i\in \mathds{N}} x_i} < 2 \cdot \varepsilon_0 $, es gilt $T \enbrace*{\sum_{i \in \mathds{N}}x_i}=y$.
	($x_i \in N$, falls $B(0,2 \varepsilon_0) \subset N$)
	\item Sei $U \subset X$ offen. Dann ist $T(U) \subset Y$ offen:
	
	Zu $x \in U$ existiert $0 \in N \subset X$ offen mit $x+ N \subset U$, also $T(x)+ T(N)\subset T(U)$. Nach (iii) existiert $0 \in Z \subset Y$ offen mit $Z \subset T(N)$,
	also $T(x)+ Z \subset T(U)$. Aber $T(x) \in T(x)+ Z$ ist offen in $Y$ und $T(U)$ ist Umgebung von $T(x)$. $x \in U$ war beliebig, also ist $T(U)$ offen. \bewende
\end{enumerate}
% subsection 38 (end)

\subsection{Corollar: Satz von der inversen Abbildung} % (fold)
\label{sub:39}
Seien $X,Y$ Banachräume, $T \in \mathcal{L}(X,Y)$ bijektiv. Dann gilt $T ^{-1} \in \mathcal{L}(Y,X)$, d.h. $T$ ist ein Homöomorphismus\index{Homöomorphismus}.
\minisec{Beweis}
Nach Satz \ref{sub:38} ist $T$ offen und somit ist $T ^{-1}$ stetig. \bewende
% subsection 39 (end)

\subsection{Corollar} % (fold)
\label{sub:310}
Sei $X$ ein Vektorraum und seien $\mathcal{T}_1, \mathcal{T}_2$ zwei Topologien auf $X$, die durch Banachraumnormen $\norm[1]{\cdot }$ und $\norm[2]{\cdot}$ induziert sind.
Falls $\mathcal{T}_1 \subset \mathcal{T}_2$, so gilt bereits $\mathcal{T}_1 = \mathcal{T}_2$.
\minisec{Beweis}
$\mathcal{T}_1 \subset \mathcal{T}_2$ $\Rightarrow $ $(X,\mathcal{T}_1) \xrightarrow{\id_X} (X,\mathcal{T}_2)$ ist stetig und bijektiv, also ist $\id_X$ nach 
\ref{sub:39} ein Homöomorphismus. \bewende
% subsection 310 (end)

\subsection{Corollar} % (fold)
\label{sub:311}
Sei $X$ ein Banachraum mit abgeschlossenen Teilräumen $X_1,X_2$, sodass $X_1 \cap X_2 = \set{0}$ und $X= X_1 + X_2$. Dann sind $X$ und $X_1 \times X_2$ als topologische
Vektorräume isomorph.
\minisec{Beweis}
Die Abbildung $X_1 \times X_2 \to X$, $(x_1,x_2) \mapsto x_1 + x_2$ ist linear und bijektiv. Sie ist stetig, wegen 
\[
	\norm[X]{x_1 + x_2} \le \norm[X]{x_1} + \norm[X]{x_2} = \norm[1,X_1 \times X_2]{(x_1,x_2)} 
\]
Mit Corollar \ref{sub:39} folgt die Behauptung. \bewende
% subsection 311 (end)
% section 3 (end)
\newpage

\section{$L^p$-Räume und der Satz von Riesz-Fischer} % (fold)
\label{sec:4}

\subsection[Erinnerung: $L^1(\mu)$ ist ein Banachraum]{Erinnerung} % (fold)
\label{sub:41}
Sei $(X,\Sigma, \mu)$ ein \Index{Maßraum}.\footnote{$X$ Menge, $\Sigma$ eine $\sigma$-Algebra, $\mu$ Maß}  $f : X \to (-\infty,\infty]$ heißt \Index{messbar}, falls 
$f ^{-1}(E) \in \Sigma$ gilt für $E \subset \mathcal{B}(\mathds{R})$ (Borelmengen\todo{Def. hinzufügen}). $s : X \to (-\infty,\infty]$ heißt \Index{einfach}, falls 
$s(X)$ endlich. Für $f : X \to [0,\infty]$ messbar setze
\[
	\int f  \, \mathd \mu := \sup \set[\int s \, \mathd \mu]{s \text{ einfach, messbar }, 0 \le s \le f} 
\] 
Für $f : X \to (-\infty,\infty]$ messbar setze $\int\! f \, \mathd \mu := \int \! f_+ \, \mathd \mu - \int \! f_- \, \mathd \mu$, falls $\int \! f_- \, \mathd \mu < \infty$.
Betrachte
\begin{align*}
	\mathcal{L}^1(\mu) := \set[f : X \to \mathds{R} \text{ messbar}]{\int \abs*{f}\,\mathd_\mu < \infty } 
\end{align*}
mit der Halbnorm $\norm[1]{f} := \int \! \abs*{f} \, \mathd_\mu $. Setze nun $\mathcal{N}^1 := \set[f : X \to \mathds{R} \text{ messbar}]{\int \abs*{f} \, \mathd_\mu =0 }$.
$L^1(\mu)$ ist ein topologischer Vektorraum (Topologie induziert durch $\norm[1]{\cdot}$) und $\mathcal{N}^1 \subset \mathcal{L}^1(\mu)$ ist ein abgeschlossener 
Untervektorraum (Übung). Definiere 
\[
	L^1(\mu) := \nicefrac{\mathcal{L}^1}{\mathcal{N}^1}
\]
$\norm[1]{\cdot }$ induziert eine Norm $\norm[1]{\cdot }$ auf $L^1(\mu)$ via  
\[
	\norm[1]{f+ \mathcal{N}^1} := \norm[1]{f} \quad , \quad f \in \mathcal{L}^1(\mu) \tag*{(Übung)}
\]
$(L^1(\mu),\norm[1]{\cdot})$ ist ein Banachraum (Ana \RM{3}, 9.23).\todo{CopyPaste aus Ana3}
% subsection 41 (end)

\subsection[Beispiele für $L^1$-Räume]{Beispiele} % (fold)
\label{sub:42}
\begin{enumerate}[(i)]
	\item $L^1(\mathds{R}^n) := L^1(\mathds{R}^n, \Lambda, \lambda)$, wobei $\Lambda$ die Lebesgue-messbaren Mengen sind und $\lambda $ das Lebesgue-Maß ist.
	\item $\ell^1(\mathds{N})= L^1 \enbrace[\big]{\mathds{N}, \mathcal{P}(\mathds{N}), \delta} = \set[f : \mathds{N} \to \mathds{R}]{\sum \abs*{f(n)} < \infty } $, 
	wobei $\delta$ das Zählmaß\footnote{siehe auch \url{http://de.wikipedia.org/wiki/Zählmaß_(Maßtheorie)}} ist.
\end{enumerate}
% subsection 42 (end)

\subsection[Definition und Proposition: Der Raum $L^p(\mu)$]{Definition und Proposition} % (fold)
\label{sub:43}
Für einen Maßraum $(X,\Sigma,\mu)$ definiere 
\begin{align*}
	\mathcal{L}^p(\mu) &:= \set[f : X \to \mathds{R} \text{ messbar}]{\int \abs*{f}^p \mathd_\mu < \infty } \\
	\mathcal{N}^p &:= \set[f : X \to \mathds{R} \text{ messbar}]{\int \abs*{f}^p \mathd_\mu = 0 } 
\end{align*}
$\mathcal{L}^p(\mu)$ und $\mathcal{N}^p$ sind Vektorräume; wir setzen 
\[
	L^p(\mu) := \nicefrac{\mathcal{L}^p(\mu)}{\mathcal{N}^p}
\]
Definiere $\norm[p]{\cdot } : \mathcal{L}^p(\mu) \to [0,\infty)$ durch $\norm[p]{f} := \enbrace*{\int \abs*{f}^p \mathd_\mu }^{1/p}$. $\norm[p]{\cdot }$ induziert eine 
Abbildung $\norm[p]{\cdot } : L^P(\mu) \to [0,\infty) $ via $\norm[p]{f + \mathcal{N}^p} := \norm[p]{f} $, für $f \in \mathcal{L}^p$.
\minisec{Beweis}
\begin{description}
	\item[$\mathcal{L}^p(\mu)$ ist Vektorraum:]  Für $f,g \in \mathcal{L}^p(\mu)$ und $\alpha, \beta \in \mathds{R}$
	\begin{align*}
		\int \abs*{\alpha \cdot f + \beta \cdot g}^p \mathd_\mu \le \int \enbrace[\big]{\abs*{\alpha \cdot f} + \abs*{\beta \cdot g}}^p \mathd_\mu
		&\le \int \enbrace[\Big]{2 \cdot \sup \set{\abs*{\alpha \cdot f}, \abs*{\beta \cdot g}  } }^p \mathd_\mu \\ 
		&\le \int \enbrace*{2^p \cdot \sup \set{\abs*{\alpha \cdot f}^p, \abs*{\beta \cdot g}^p  \rule{0cm}{1em}} } \mathd_\mu \\ 
		&\le \int 2^p \cdot  \enbrace[\big]{\abs*{\alpha \cdot f}^p + \abs*{\beta \cdot g}^p} \mathd_\mu < \infty 
	\end{align*}
	\item[$\mathcal{N}^p$ ist Vektorraum:] Übung.
	\item[{$\norm[p]{\cdot}$ wohldefiniert auf $L^p(\mu)$:}] Übung. \bewende
\end{description}
% subsection 43 (end)

\subsection[Proposition: Abschätzung für $\alpha, \beta \ge 0$ und $p,q >1$ mit $\frac{1}{p} + \frac{1}{q}=1$]{Proposition} % (fold)
\label{sub:44}
Für $\alpha, \beta \ge 0$, $p,q >1$ mit $\frac{1}{p} + \frac{1}{q}=1$. Dann gilt 
\begin{equation*}
	\alpha \cdot \beta \le \frac{\alpha^p}{p} + \frac{\beta^q}{q}  \label{eq:44}\tag{\#}
\end{equation*}
\minisec{Beweis}
\begin{align*}
	\frac{1}{p} + \frac{1}{q}=1 \iff p+q= pq \iff (p-1)(q-1)= 1 \iff p (q-1) = q \iff q (p-1) =p
\end{align*}
Falls $\beta = \alpha^{p-1}$, so gilt $\alpha \cdot \beta = \alpha^p$ und 
\[
	\frac{\alpha^p}{p} + \frac{\beta^q}{q} = \frac{\alpha^p}{p} + \frac{(\alpha^{p-1})^q}{q} = \frac{\alpha^p}{p} + \frac{\alpha^p}{q} = \alpha^p = \alpha \cdot \beta   
\]
Also gilt Gleichheit in \eqref{eq:44}. Für festes $\alpha \ge 0$ betrachte $f_\alpha : [0,\infty) \to [0,\infty)$ gegeben durch
$f_\alpha(\beta) := \frac{\alpha^p}{p} + \frac{\beta^q}{q} - \alpha \cdot \beta$. Dann ist $f_\alpha$ stetig auf $[0,\infty)$ und differenzierbar auf $(0,\infty)$. Es gilt
$f_\alpha'(\beta) = \beta^{q-1} -\alpha$ und weiter
\[
	f_\alpha'(\beta)=0 \iff \beta^{q-1} = \alpha \iff \beta= \alpha^{p-1}
\]
$f_\alpha''(\alpha^{p-1}) \ge 0$, also hat $f_\alpha$ bei $\beta = \alpha^{p-1}$ ein lokales Minimum. Es gilt $f_\alpha(\alpha^{p-1})= 0 \le f_\alpha(0)$, also
hat $f_\alpha$ bei $\beta=\alpha^{p-1}$ ein globales Minimum. $\Rightarrow $ \eqref{eq:44}. \bewende
% subsection 44 (end)

\subsection{Satz: Höldersche Ungleichung} % (fold)
\label{sub:45}
Sei $(X,\Sigma,\mu)$ ein Maßraum. Seien $p,q \ge 1$ mit $\frac{1}{p} +\frac{1}{q}=1 $, $f \in \mathcal{L}^p(\mu)$, $g \in \mathcal{L}^q(\mu)$. Dann gilt 
$f \cdot g \in \mathcal{L}^1{\mu}$ und
\[
	\norm[1]{f \cdot g} \le \norm[p]{f} \cdot \norm[q]{g}   
\]
Die entsprechenden Aussagen gelten dann auch für $L^p(\mu)$, $L^q(\mu)$, $L^1(\mu)$.
\minisec{Beweis}
Klar, falls $\norm[p]{f} =0 $ oder $\norm[q]{g} =0$ (Warum?). Sei also $\norm[p]{f}, \norm[q]{g} \not= 0$. Nach Proposition \ref{sub:44} gilt für $x \in X$
\[
	0 \le \underbrace{\underbrace{\frac{\abs*{f(x)} }{\norm[p]{f}}}_{=\alpha} \cdot \underbrace{\frac{\abs*{g(x)} }{\norm[q]{g} }}_{=\beta} }_{\text{messbar}}
	\le \underbrace{\frac{1}{p} \cdot \frac{\abs*{f(x)}^p }{\norm[p]{f}^p } + \frac{1}{q} \cdot \frac{\abs*{g(x)}^q }{\norm[q]{g}^q }}_{\text{integrierbar}}
\]
$\abs*{f}^p , \abs*{g}^q  \in \mathcal{L}^1(\mu)$, also $\abs*{ f \cdot g } \in \mathcal{L}^1(\mu)$. Es gilt
\begin{align*}
	\frac{\int \abs*{f \cdot g} \mathd_\mu}{\norm[p]{f} \cdot \norm[q]{g}} \le \frac{1}{p} \cdot \frac{\int \abs*{f}^p \mathd_\mu }{\norm[p]{f}^p } +
	\frac{1}{q} \cdot \frac{\int \abs*{g}^q \mathd_\mu }{\norm[q]{g}^q} = \frac{1}{p} + \frac{1}{q} = 1       
\end{align*}
$\Rightarrow \inf \abs*{f \cdot g} \mathd_\mu  \le \norm[p]{f}  \cdot \norm[q]{g}$. \bewende
% subsection 45 (end)

\subsection{Satz: Minkowskische Ungleichung} % (fold)
\label{sub:46}
Für $p>1$, $f,g \in \mathcal{L}^p(\mu)$ (bzw. $L^p(\mu)$) gilt 
\[
	\norm[p]{f +g} \le \norm[p]{f} + \norm[p]{g}   
\]
\minisec{Beweis}
Sei $\norm[p]{f+g} >0 $ (sonst trivial). Es gilt 
\begin{align*}
	\abs*{f +g}^p = \abs*{f+g}^{p-1} \cdot \abs*{f+g} \le \abs*{f+g}^{p-1} \enbrace[\big]{\abs*{f} + \abs*{g}}     
\end{align*}
Sei $q>1$ mit $\frac{1}{p} + \frac{1}{q}=1$. Es gilt $\abs*{f+g}^{p-1} \in \mathcal{L}^q(\mu)$, da $\abs*{f+g}^{(p-1)q} = \abs*{f+g}^p \in \mathcal{L}^1(\mu)$. Nach Höder
\ref{sub:45} gilt
\begin{align*}
	\int \abs*{f+g}^p \mathd_\mu \le \int \abs*{f+g}^{p-1} \abs*{f} \mathd_\mu + \int \abs*{f+g}^{p-1} \abs*{g}\mathd_\mu &\le \norm[p]{f} \cdot \norm[q]{\abs*{f+g}^{p-1}}
	+ \norm[p]{g} \cdot \norm[q]{\abs*{f+g}^{p-1} } \\        
	&= \enbrace*{\norm[p]{f} + \norm[p]{g}} \cdot \enbrace*{\int \abs*{f+g}^{(p-1)q} \mathd_\mu }^{\frac{1}{q}} \\
	&= \enbrace*{\norm[p]{f} + \norm[p]{g}} \cdot \enbrace*{\int \abs*{f+g}^{p} \mathd_\mu }^{q} \\
\end{align*}
Also gilt 
\[
	\enbrace*{\int \abs*{f+g}^p \mathd_\mu}^{1- \frac{1}{q}} = \enbrace*{\int \abs*{f+g}^{p} \mathd_\mu }^{\frac{1}{p}} \le \norm[p]{f} + \norm[p]{g} \bewende    
\]
% subsection 46 (end)
% section 4 (end)



\cleardoubleoddemptypage
\appendix
\section{Anhang} % (fold)
\label{sec:anhang}

\subsection{Vierecksungleichung} % (fold)
\label{sub:vier_ungl}
Sei $(X,d)$ ein metrischer Raum und $x,y,u,v \in X$. Dann gilt
\[
	\abs[\big]{d(x,y) - d(u,v)} \le d(x,u) + d(y,v) 
\]
\minisec{Beweis}
Einerseits gilt nach der Dreiecksungleichung
\[
	d(x,y) \le d(x,u) + d(u,v) + d(v,y) \quad \Longrightarrow \quad d(x,y) - d(u,v) \le d(x,u) + d(y,v)
\]
Andererseits aber auch
\[
	d(u,v) \le d(u,x) + d(x,y) + d(y,v) \quad \Longrightarrow \quad d(u,v) - d(x,y) \le d(x,u) + d(y,v)
\]
Insgesamt folgt also die Behauptung. \bewende
% subsection vier_ungl (end)

Tatsache für Übungen: $X,Y$ Banachräume, $T \in \mathcal{L}(X,Y)$ surjektiv $\Rightarrow T$ ist offen.

Corollar: $T \in \mathcal{L}(X,Y)$ bijektiv $\Rightarrow $ $T ^{-1}$ stetig.

% section anhang (end)
\cleardoubleoddemptypage
\pagenumbering{Alph}
\setcounter{page}{1}
\printindex
\listoffigures
\todototoc
\listoftodos[Todo's und andere Baustellen]
\end{document}
